% Options for packages loaded elsewhere
\PassOptionsToPackage{unicode}{hyperref}
\PassOptionsToPackage{hyphens}{url}
\documentclass[
]{article}
\usepackage{xcolor}
\usepackage[margin=1in]{geometry}
\usepackage{amsmath,amssymb}
\setcounter{secnumdepth}{-\maxdimen} % remove section numbering
\usepackage{iftex}
\ifPDFTeX
  \usepackage[T1]{fontenc}
  \usepackage[utf8]{inputenc}
  \usepackage{textcomp} % provide euro and other symbols
\else % if luatex or xetex
  \usepackage{unicode-math} % this also loads fontspec
  \defaultfontfeatures{Scale=MatchLowercase}
  \defaultfontfeatures[\rmfamily]{Ligatures=TeX,Scale=1}
\fi
\usepackage{lmodern}
\ifPDFTeX\else
  % xetex/luatex font selection
  \setmainfont[Ligatures=TeX]{Cambria}
  \setmonofont[]{Consolas}
\fi
% Use upquote if available, for straight quotes in verbatim environments
\IfFileExists{upquote.sty}{\usepackage{upquote}}{}
\IfFileExists{microtype.sty}{% use microtype if available
  \usepackage[]{microtype}
  \UseMicrotypeSet[protrusion]{basicmath} % disable protrusion for tt fonts
}{}
\makeatletter
\@ifundefined{KOMAClassName}{% if non-KOMA class
  \IfFileExists{parskip.sty}{%
    \usepackage{parskip}
  }{% else
    \setlength{\parindent}{0pt}
    \setlength{\parskip}{6pt plus 2pt minus 1pt}}
}{% if KOMA class
  \KOMAoptions{parskip=half}}
\makeatother
\usepackage{longtable,booktabs,array}
\newcounter{none} % for unnumbered tables
\usepackage{calc} % for calculating minipage widths
% Correct order of tables after \paragraph or \subparagraph
\usepackage{etoolbox}
\makeatletter
\patchcmd\longtable{\par}{\if@noskipsec\mbox{}\fi\par}{}{}
\makeatother
% Allow footnotes in longtable head/foot
\IfFileExists{footnotehyper.sty}{\usepackage{footnotehyper}}{\usepackage{footnote}}
\makesavenoteenv{longtable}
\usepackage{graphicx}
\makeatletter
\newsavebox\pandoc@box
\newcommand*\pandocbounded[1]{% scales image to fit in text height/width
  \sbox\pandoc@box{#1}%
  \Gscale@div\@tempa{\textheight}{\dimexpr\ht\pandoc@box+\dp\pandoc@box\relax}%
  \Gscale@div\@tempb{\linewidth}{\wd\pandoc@box}%
  \ifdim\@tempb\p@<\@tempa\p@\let\@tempa\@tempb\fi% select the smaller of both
  \ifdim\@tempa\p@<\p@\scalebox{\@tempa}{\usebox\pandoc@box}%
  \else\usebox{\pandoc@box}%
  \fi%
}
% Set default figure placement to htbp
\def\fps@figure{htbp}
\makeatother
\setlength{\emergencystretch}{3em} % prevent overfull lines
\providecommand{\tightlist}{%
  \setlength{\itemsep}{0pt}\setlength{\parskip}{0pt}}
\usepackage{xurl}
\usepackage{bookmark}
\IfFileExists{xurl.sty}{\usepackage{xurl}}{} % add URL line breaks if available
\urlstyle{same}
\hypersetup{
  pdftitle={Frontier and Localhost: How Production AI Learns Outside the Weights},
  pdfauthor={Mikhail L. Arbuzov; Alexey A. Shvets; Sisong Bei},
  hidelinks,
  pdfcreator={LaTeX via pandoc}}

\title{Frontier and Localhost: How Production AI Learns Outside the
Weights}
\author{Mikhail L. Arbuzov \and Alexey A. Shvets \and Sisong Bei}
\date{}

\begin{document}
\maketitle

\textbf{Abstract}

Production LLM systems increasingly adapt through scaffold artifacts
outside the model weights: prompts, rules, memories, skills, tools,
routing graphs, eval suites, and governance pipelines that update from
deployment feedback on a cadence the weights cannot match. Yet this
layer is still mostly maintained as \emph{patchwork}: failures trigger
local fixes; fixes accumulate as brittle prompts, stale memories, ad hoc
tool bindings, and unprincipled workflow exceptions. The field has built
an external adaptation layer for LLMs without building the optimizer
that should govern it. Surveying approximately 130 production and
research systems from 2023--2026, we formalise the disciplined
alternative as \textbf{artifact-layer descent}: a discrete, gated,
deployment-time update process in which recurring residuals generate
candidate scaffold deltas, gates estimate finite directional loss
change, accepted deltas persist, and promotion across contexts is
bounded by evidence radius. The resulting two-loop architecture (local
update; cross-context promotion under a gate) is consistent with
convergent patterns across IDE plugins, vertical-bundle vendors, and
self-evolving research agents. Thirteen research systems plus two
industry exemplars satisfy a strict composite-two-loop criterion;
placing them on a 3×3 matrix of gate types leaves three populated cells
and three empty cells. The consequential empty cell is enterprise
full-auto: automated local evolution combined with governed cross-tenant
promotion, versioned lineage with RBAC, and rollback --- in ML language,
a DP-FedAvg-style cross-tenant aggregator with an automated eval gate at
the scaffold layer, in the limit where artifacts are embedded as
mergeable update representations. The contribution is not the claim that
prompts are weights; it is that production AI now has an external
adaptation surface that needs optimization, regularization, versioning,
rollback, and governance, and that the path from today's patchwork to
that discipline is a tractable engineering programme. A
practitioner-facing companion paper will handle reference architectures,
maturity ladders, design principles, and the full failure taxonomy.

\begin{center}\rule{0.5\linewidth}{0.5pt}\end{center}

\subsection{1. The gradient moved outside the
model}\label{the-gradient-moved-outside-the-model}

The standard mental model of an LLM system is the model. Pre-training
fixes the weights; fine-tuning or RLHF nudges them; deployment wraps a
thin prompt-plus-tool layer around the result and ships. Errors are
weight errors; reliability is a weight-training problem.

Production has stopped behaving this way. In 2025--2026, frontier models
update on cycles measured in months, but the systems built on top of
them update continuously --- in markdown files, persistent memories,
retrieved skills, registered tools, evaluation suites, agent topologies,
version-controlled prompt repositories, and PR-gated governance. A
growing share of deployment-time adaptation now occurs outside the
weights, in a fast-changing \emph{localhost} scaffold wrapped around a
slow-changing frontier. The section title is a metaphor --- the gradient
inside the model still happens during training, and the scaffold layer
is not a gradient in the calculus sense. The substantive claim is the
milder one: production AI has accidentally created an external learning
surface, and right now it is being operated as patchwork rather than as
a disciplined optimization process. What this paper formalises is the
disciplined limit case --- a discrete, gated, deployment-time update
process that adapts a frozen frontier to the recurring failure topology
of a specific patch --- and what the survey finds is that almost no
deployed system has yet built the optimizer for that surface.

This paper argues that scaffold evolution can be modeled as learning in
a precise sense: scaffold artifacts are updated from failure signal, the
updated artifacts reduce future failure probability on a local patch,
and the update process --- observe failure, generate candidate delta,
gate it, persist it, optionally promote it --- has the structure of
discrete, gated, stochastic local optimization. The substrate is
symbolic and human-readable rather than continuous and differentiable,
and the gate stands in for the gradient; the model is best read as
\emph{artifact-layer optimization}, with descent-like behavior under the
assumptions catalogued in §3.2. Where those assumptions fail (no
measurable patch loss, no credit-assignment process, no persistent
artifact distinct from optimiser state) the framing should be read as
modeling language rather than a theorem.

Read together, these artifacts form a single external adaptation surface
that current production practice still treats as scattered engineering.
Naming the surface makes the optimization discipline it needs --- and
visibly lacks --- visible. A Claude Skill is not ``just a markdown
file.'' It is a learned procedural weight kept outside the model. A
Cursor Rule is not ``just an instruction.'' It is a local parameter for
one project's distribution shift. A RAG connector is not ``just
retrieval.'' It is patch-specific memory access. A tool registry is not
``just integration.'' It is capability provisioning along an axis the
base model cannot reach. A PR review against a shared rule repository is
not ``just governance.'' It is a promotion gate on a federated update.
An eval suite that blocks merge is not ``just testing.'' It is a
validation-loss check before promotion. A cross-tenant artifact
repository is not ``just sharing.'' It is \emph{federated artifact
promotion} across distributed shards --- scope expansion under a gate,
rather than arithmetic averaging of numeric updates (we sharpen the
distinction in §3.2 and §9). The mechanisms were already in production;
what is new is reading them as a single external adaptation surface
rather than as scattered engineering choices, and asking what
optimization discipline that surface needs.

We survey approximately 130 systems across six substrates (instructions,
skills, memory, tools, orchestration, governance), score each on a 0--5
patch-plasticity rubric, and analyse the thirteen research systems (plus
two industry exemplars) that satisfy a strict \emph{composite two-loop}
criterion: an inner loop that updates a persisted artifact from session
feedback, an outer loop that promotes the artifact across scopes under
explicit versioning. The two-loop design space is a 3×3 matrix on gate
type --- \emph{Auto-gated} when the commit-or-promote decision is
operationalised by an algorithmic criterion, \emph{Human-gated} when it
is reviewer judgment. The matrix has three populated cells; three are
empty for distinct structural reasons. The fully automated, governed,
cross-tenant cell has no surveyed occupant. That empty corner is the
paper's open-question headline and a concrete Paper-4 target.

The architectural review of how to actually build such systems --- five
reference patterns, a five-level maturity model, a substrate-selection
rubric, three migration paths, the full nine-category failure taxonomy,
and eight design principles --- is deferred to follow-on practitioner
work. This paper is the survey and the framework. The single-line thesis
is this: \textbf{the field has built an external adaptation layer for
LLMs, but not the optimizer that should govern it.} The rest of the
paper develops that observation into a survey of where production sits
on a \emph{patchwork → random-search → descent} gradient, a
formalisation of the disciplined limit case under stated assumptions,
and a named empty corner where the optimizer most plainly does not yet
exist.

\subsection{2. Why patch errors must live outside the
weights}\label{why-patch-errors-must-live-outside-the-weights}

The scaffold is not merely where production systems happen to store
patches. For high-churn, tenant-specific, auditable adaptation, the
scaffold is \emph{often the practical substrate} --- not because no
alternative exists, but because under deployment-time constraints
(release cadence, tenant isolation, inspectability, rollback) it is
cheap, locally scoped, reversible, and updatable at human time-scales
while the weights are not. Frontier training and production reliability
optimise for objectives that are in tension, and that tension pushes
most patch-specific adaptation outside the model. Per-tenant fine-tunes,
LoRA adapters, distilled small models, and learned routing can in
principle absorb some patch residuals; the survey finds production has
consequently located the bulk of patch-residual adaptation in the
scaffold, but whether this remains true under different deployment
constraints is an empirical question rather than a structural necessity.

\textbf{Frontier training optimises for breadth.} The weights are tuned
for cross-patch generalisation: they smooth over local conventions,
contradictory tenant preferences, and rare workflow-specific failures in
order to preserve broad competence across millions of unseen
deployments. The smoothing is not incidental. It is the training
objective. A model that overfits to one tenant's naming conventions or
one team's preferred error-handling style loses transfer to the next
tenant; a model that encodes every patch's idiosyncratic eval
expectation loses the general-purpose competence that made it worth
deploying in the first place. The smoothing dynamics have been measured
directly: RLHF reduces output diversity and homogenises preferred styles
across users (Kirk et al., 2024; Padmakumar \& He, 2023), and
instruction-tuning data distributions favour patterns that generalise
across the annotator pool rather than patterns specific to any one
deployment. The frontier model is therefore \emph{trained to be unable}
to encode a single tenant's idiosyncrasies as preferred behaviour;
absorbing that idiosyncrasy without weight retraining is exactly what
the scaffold layer is for.

\textbf{Production reliability is achieved in depth.} Inside a single
deployment patch, reliability is determined by exactly the stuff global
training had to smooth away: local naming conventions, team-specific
workflows, domain-specific failure patterns, project-specific schemas,
idiosyncratic tool chains, customer-specific risk tolerances, recurring
local mistakes, institutional preferences, hidden evaluator
expectations. The per-patch quantities from our previous work ---
failure-mode catalogue size \(|C_D|\), hard-fraction \(\beta_D\),
mode-discovery rate \(\sigma_D\) --- are local by construction, and a
library calibrated to one patch under-covers the next.

\textbf{Encoding every patch in the weights is structurally infeasible.}
Three constraints rule out the obvious approach. \emph{Economic:}
millions of patches multiplied by billions of parameters times the cost
of per-patch fine-tuning is not a viable training-compute budget.
\emph{Release-cycle:} the patch wants daily updates; weight releases
ship quarterly at best. \emph{Cross-patch interference:} overfitting the
weights to one tenant's preferences degrades behaviour for adjacent
tenants whose preferences contradict. The weights are tuned for
cross-patch invariants by construction; that is what they are
\emph{for}.

The localhost scaffold solves the conflict. It is the substrate that
uniquely satisfies all of:

\begin{itemize}
\tightlist
\item
  overfit aggressively to one tenant without damaging another;
\item
  encode local conventions without polluting global behaviour;
\item
  update daily instead of waiting for release cycles;
\item
  be inspected, versioned, audited, and rolled back;
\item
  attach tools, memories, rules, workflows, and evals that are
  meaningless outside the patch;
\item
  specialise inside the patch while the model remains broadly competent.
\end{itemize}

Each scaffold substrate targets a different patch-residual class.
Instructions and tools address \emph{capability} residuals --- the model
could not do this on its own (PAL for arithmetic, RAG for retrieval,
constrained decoders for format). Skills and memory address
\emph{catalogue} residuals --- the same patch-specific failure mode
keeps recurring and now has a patch. Orchestration addresses
\emph{hard-fraction} residuals --- the per-step decision was too hard,
so specialised sub-agents bring it back into capability range.
Governance addresses \emph{propagation} residuals --- a correction
proven in one deployment is wasted unless it can be promoted to others
under audit.

Our previous work (Arbuzov, Shvets \& Bei, 2025; \emph{Architecture of
Errors}, 2026) argued that LLM errors concentrate sparsely at key tokens
and cluster into a finite catalogue of recurring modes whose size grows
polylogarithmically in observed failures. The framework left a
\emph{placement} question --- where do the interventions live? The
answer follows from the tension above. Under deployment-time constraints
(release cadence, tenant isolation, inspectability, rollback), the
scaffold is the substrate that satisfies all four; the survey finds that
production systems consequently locate the bulk of patch-residual
interventions there. The frontier model generalises; the localhost
scaffold specialises; reliability comes from governing that
specialisation.

\subsection{2.5 Patchwork, random artifact search, and artifact-layer
descent}\label{patchwork-random-artifact-search-and-artifact-layer-descent}

Three stages of scaffold maintenance practice recur across the corpus.
They differ in how candidate updates are produced, in whether
failure-mode discovery and credit assignment are explicit, and in
whether promotion is bounded by evidence radius.

\textbf{Stage 1 --- Patchwork.} Humans add prompts, rules, tools,
memories, and routing logic after failures. Each fix is proposed by
intuition, committed with weak evidence, and accumulated without
pruning. The dominant failure mode is bloat: scaffolds grow until
lost-in-the-middle (§10) erases the marginal benefit of new additions.
Most production deployments observed in the survey sit here.

\textbf{Stage 2 --- Random artifact search.} Many scaffold variants are
tried --- A/B'd rubrics, swapped memory schemas, regenerated tool
bundles --- but without stable credit assignment or principled
promotion. Improvements happen locally; transfer is unreliable;
overfitting to the in-distribution patch is the dominant failure mode.
The governance-first cluster of §7 sits here when the eval pipeline is
automated but the candidate-generation pipeline is not.

\textbf{Stage 3 --- Artifact-layer descent.} Candidate deltas are tied
to recurring failure modes, scored by an estimator of finite directional
loss change, regularized by complexity/risk penalties, accepted only
when they clear a gate, versioned, prunable, and promoted only within
evidence radius. The composite-two-loop systems of §8 instantiate Stage
3 for \emph{within-organisation} scope; the empty corner of §9 is Stage
3 for \emph{cross-tenant} scope. The math of §3.2 applies to Stage 3 and
only suggestively to Stages 1--2.

{\def\LTcaptype{none} % do not increment counter
\begin{longtable}[]{@{}
  >{\raggedright\arraybackslash}p{(\linewidth - 8\tabcolsep) * \real{0.2000}}
  >{\raggedright\arraybackslash}p{(\linewidth - 8\tabcolsep) * \real{0.2000}}
  >{\raggedright\arraybackslash}p{(\linewidth - 8\tabcolsep) * \real{0.2000}}
  >{\raggedright\arraybackslash}p{(\linewidth - 8\tabcolsep) * \real{0.2000}}
  >{\raggedright\arraybackslash}p{(\linewidth - 8\tabcolsep) * \real{0.2000}}@{}}
\toprule\noalign{}
\begin{minipage}[b]{\linewidth}\raggedright
Stage
\end{minipage} & \begin{minipage}[b]{\linewidth}\raggedright
Update mechanism
\end{minipage} & \begin{minipage}[b]{\linewidth}\raggedright
Credit assignment
\end{minipage} & \begin{minipage}[b]{\linewidth}\raggedright
Promotion discipline
\end{minipage} & \begin{minipage}[b]{\linewidth}\raggedright
Dominant failure mode
\end{minipage} \\
\midrule\noalign{}
\endhead
\bottomrule\noalign{}
\endlastfoot
\textbf{1. Patchwork} & Intuition-driven local edits & Implicit / ad hoc
& None & Bloat \\
\textbf{2. Random artifact search} & Many variants tried under eval &
Weak; eval is the only signal & Eval-gated CI but no
failure-to-coordinate tracing & Overfitting to in-distribution patch \\
\textbf{3. Artifact-layer descent} & Residual → coordinate → delta &
Explicit per failure mode & Evidence-radius-bounded; reversible & (the
discipline this paper formalises) \\
\end{longtable}
}

The paper's central claim is therefore not ``modern AI systems are doing
stochastic descent'' --- they mostly are not. It is that \emph{they can
become so} once the patchwork loop is instrumented with measurable patch
loss, residual capture, failure-to-coordinate credit assignment,
candidate delta synthesis, and a gate that estimates directional loss
change before persistence or promotion. The survey of §§5--8 asks which
deployed systems have closed which of those five steps; the empty corner
of §9 names the closure that no public system has yet achieved.

\subsection{3. Artifact-layer descent}\label{artifact-layer-descent}

The remaining question is whether ``scaffold evolution'' is learning in
any technical sense or just engineering practice. The answer is: it
\emph{becomes} learning in a precise technical sense only in the Stage-3
limit of §2.5 --- only when the patchwork loop is instrumented with
measurable patch loss, residual capture, failure-to-coordinate credit
assignment, candidate delta synthesis, and a gate that estimates
directional loss change before persistence or promotion. Current
practice across the surveyed corpus mostly resembles unguided search
over scaffold configurations (Stage 1 or Stage 2); the formalism below
describes the Stage-3 limit, and the survey of §§5--8 asks which
deployed systems have closed which of the five required steps.

Let \(D\) be a deployment patch --- a specific application, user, or
task distribution. \emph{Patch} is used throughout this paper with a
defined scope hierarchy: USER ⊂ PROJECT ⊂ STACK ⊂ TENANT ⊂ CORE. A USER
patch is one human's account or fork; a PROJECT is one repository,
workflow, or product workspace; a STACK is a curated bundle of projects
under one team; a TENANT is an organisation or customer; CORE is the
cross-tenant universal scope. Each level induces its own loss
\(L_{D_{\text{level}}}\) (the average over the patches inside that
scope). When the text says ``cross-tenant promotion'' or
``cross-organisation federation,'' it specifically means moving
artifacts up to the CORE or cross-TENANT level; ``patch-specific
adaptation'' without qualifier means descent on a USER or PROJECT-level
\(L_D\). Let \(M\) be the frozen frontier model, accessed via inference
but not updated. Let \(S_t\) be the scaffold at time \(t\): the layered
collection \(\{S_t^{(1)}, \ldots, S_t^{(6)}\}\) of instruction set,
skill library, memory store, tool registry, orchestration graph, and
governance configuration. Let \(\pi(M, S_t, x)\) denote the system's
behaviour on input \(x\) --- the model running on the
model-plus-scaffold composite. Let
\(L_D(S_t) = \mathbb{E}_{x \sim D}[\ell(\pi(M, S_t, x))]\) be the
expected residual failure rate on the patch.

A failure observation produces a candidate delta \(\Delta S\): an
instruction edit, a new skill, a memory correction, a tool
re-registration, a topology change, or a governance-rule update. A gate
\(G_t\) accepts \(\Delta S\) if the estimated patch-loss change clears
an acceptance threshold \emph{and} risk stays within tolerance:

\[
S_{t+1} = \begin{cases} S_t \oplus \Delta S, & G_t(\Delta S) = 1, \\ S_t, & G_t(\Delta S) = 0, \end{cases}
\]

where \(\oplus\) denotes structured application of an admissible edit
(an instruction rewrite, a skill insertion, a memory correction, etc.)
to the current scaffold.

Different surveyed systems instantiate \(G\) differently. NanoResearch's
gate is a semantic-merge by the orchestrator (algorithmic). SkillRL's
gate is an RL threshold at task-category accuracy below 0.4
(algorithmic). AlphaEvolve's gate is a cascade evaluator with MAP-Elites
archive selection (algorithmic). SkillForge's gate is an LLM-judge with
\textgreater90\% human agreement \emph{plus} a human support engineer
(hybrid algorithmic + reviewer). Anthropic's skills repository uses
reviewer judgment on a PR-by-PR basis with no algorithmic threshold
(purely reviewer-judgment). DSPy's gate is gradient-like over
prompt-template space without persistent commitment. Reflexion's gate is
essentially absent: any reflection becomes a memory entry.

These gates serve related but not identical roles. Eval-gated,
LLM-judge-gated, RL-threshold-gated, and reviewer-gated systems all
estimate a directional loss change on the candidate edit and accept iff
that estimate clears a margin --- they instantiate the per-edit gate of
§3.2 directly. Semantic-merge and MAP-Elites archive-selection gates are
better interpreted as \emph{proposal-pruning} steps that precede the
per-edit gate: they bound which \(z_t\) enter consideration rather than
directly estimating \(\widehat{\Delta} L_D\). Appendix C maps each
surveyed gate type to its slot in the §3.2 formalism. Across the
spectrum, each system is asking, by some method, whether \(\Delta S\)
moves \(S_t\) toward lower \(L_D\). The substrate is discrete: scaffold
edits are typically textual, procedural, or symbolic, not differentiable
in the backpropagation sense. The gate is mediated: humans, evals,
LLM-judges, or surrogate verifiers stand in for the gradient. The update
is governed: not every candidate \(\Delta S\) propagates; some are
rejected, some are quarantined to a single user's fork, some climb a
layered hierarchy from local to shared scope before they take.

We call this \emph{artifact-layer descent}. In a differentiable
relaxation --- where scaffold edits are represented as coordinates in an
artifact basis, proposals are generated from estimated negative-gradient
directions over that basis, and the gate is a differentiable
validation-loss estimator --- artifact-layer descent reduces to a form
of stochastic gradient-like local search. In ordinary deployments, where
edits are discrete, proposals are LLM- or template-generated rather than
gradient-driven, and gates are eval-, judge-, or reviewer-mediated, the
process is a discrete, gated stochastic approximation to descent. The
two are continuous with one another in modeling style; the difference is
one of substrate, proposal mechanism, and gate signal-to-noise.

Two loops follow naturally. \textbf{Loop 1} is within-context: a session
feedback signal triggers an update to a persisted local artifact.
\textbf{Loop 2} is cross-context: a local artifact is promoted to a
shared scope (cross-project, cross-user, cross-organisation) under
explicit versioning. Loop 1 is the local descent step on a per-patch
loss; Loop 2 is \emph{federated artifact promotion or aggregation} ---
scope expansion under a gate, not arithmetic averaging of numeric
updates. FedAvg is a useful analogy only in the limit where artifacts
have been embedded into mergeable update representations and the
aggregation is parameter-space averaging; the DP-FedAvg framing in §9's
empty-corner target assumes that limit. Surveyed Loop-2 mechanisms are
selection, merge, curation, review, and distribution; none of the
systems in the survey performs parameter-space averaging in any
technical sense. Each loop has a gate; gates are \textbf{Auto-gated}
when the decision is an operationalised algorithmic criterion (RL
threshold, eval score, LLM-judge with quantified agreement, surrogate
verifier, benchmark cascade), \textbf{Human-gated} when the decision is
reviewer judgment (automation often precedes the gate but the gate
itself is reviewer discretion), and \textbf{None} when the loop is
absent.

These five terms --- Loop 1, Loop 2, Auto-gated, Human-gated, None ---
anchor the rest of the paper.

The framing positions the work against several adjacent literatures.
CoALA (Sumers et al., 2023) proposes a static cognitive architecture of
modular memory and structured action; we update CoALA's snapshot into a
dynamic deployment architecture with governance. The self-evolution
surveys (Gao et al.~2026; Fang et al.~2025) treat agents as learners
that update weights, representations, or optimiser-selected artifacts
--- they locate the gradient inside the model. The memory-survey of Du
(2026) and the skills-survey of Zhou et al.~(2026) cover individual
substrates with version-and-promotion machinery that mirrors S3 and S2
here but stops at the substrate boundary. ACE (2025) treats context as a
single evolving playbook. Each adjacent work touches a piece of the
picture; none develops the cross-substrate artifact-layer descent
framing or the design space it implies.

\subsubsection{3.1 The mechanism stack: from patchwork to
descent}\label{the-mechanism-stack-from-patchwork-to-descent}

The abstract update rule \(S_{t+1} = S_t + G(\Delta S)\) expands into
six operations that every patch-plastic system must perform --- whether
explicitly engineered or accidentally emergent.

\begin{enumerate}
\def\labelenumi{\arabic{enumi}.}
\tightlist
\item
  \textbf{Residual capture.} Log failures, corrections, rejected
  outputs, user edits, eval regressions, tool-call failures. The Loop-1
  signal source. Without it, \(\Delta S\) has no provenance.
\item
  \textbf{Mode discovery.} Cluster repeated failures into patch-specific
  error modes. The Paper-2 catalogue, instantiated locally per patch.
\item
  \textbf{Delta synthesis.} Convert a recurring mode into a candidate
  scaffold edit: an instruction rewrite, a new skill, a memory
  correction, a tool binding, an orchestration change, an eval case.
  This is the production of \(\Delta S\).
\item
  \textbf{Delta validation.} Estimate whether the edit reduces \(L_D\)
  without collateral damage on adjacent inputs. This is the gate \(G\)
  --- Auto-gated when the check is an algorithmic criterion (RL
  threshold, eval score, LLM-judge with quantified agreement),
  Human-gated when it is reviewer judgment.
\item
  \textbf{Promotion control.} Decide whether the accepted \(\Delta S\)
  stays in the local fork or climbs the scope hierarchy --- project,
  stack, organisation, cross-tenant. This is Loop 2.
\item
  \textbf{Drift and bloat control.} Expire stale rules, prune
  conflicting memories, roll back regressed skills, detect over-fitting
  to one patch. Without this layer, accumulated \(\Delta S\) produces
  the scaffold-level analogue of plasticity loss.
\end{enumerate}

Each operation is the locus of a separate engineering discipline that
does not yet exist in coordinated form. Surveyed systems implement the
operations partially and unevenly: research full-auto systems excel at
(3) and (4), production governance leaders at (5), industry hybrids at
most of (1)--(5) but rarely (6). Implementation patterns for each of the
six are deferred to follow-on work; this paper's purpose is to make the
discipline visible.

\subsubsection{3.2 The mechanism made formal: scaffold updates as
discrete stochastic
descent}\label{the-mechanism-made-formal-scaffold-updates-as-discrete-stochastic-descent}

\textbf{Assumptions and scope of the formalization.} The descent
argument below holds when: (i) \(\theta_M\) is fixed during the
adaptation window; (ii) scaffold edits are persistent, inspectable
artifacts distinct from transient context; (iii) the patch admits a
measurable operational loss; (iv) the proposal process can attribute
failures to scaffold coordinates with non-zero probability; (v) gates
are imperfect estimators of finite directional loss change rather than
oracle access to \(\Delta L_D\); (vi) Loop-2 promotion is \emph{scope
expansion under a gate}, not arithmetic averaging of comparable updates.
Survey claims of the form ``no surveyed system has X'' are bounded by
the audited public corpus and the survey cutoff of May 2026. Where the
formalism is metaphor rather than mechanism (a deployment lacks
measurable loss, or proposal generation cannot assign credit to a
coordinate), the descent picture should be read as modeling language
rather than a theorem about that deployment.

The mechanism stack above makes artifact-layer descent operational, but
the word ``descent'' still needs to earn its keep. A skeptical reading
is available: scaffold updates are just engineering edits, not
gradients; PR review is not backpropagation; an LLM-judge is not a
derivative; a markdown rule is not a differentiable parameter. All true.
But they do not defeat the claim. They locate the claim in the correct
optimization class.

Scaffold evolution is not backpropagation through continuous weights. It
is discrete, gated, patch-local stochastic descent over a structured
artifact space.

Let the frontier model have fixed parameters \(\theta_M\). Let the
scaffold be a structured parameter object

\[\theta_s \in \Theta_s = \Theta_1 \times \Theta_2 \times \Theta_3 \times \Theta_4 \times \Theta_5 \times \Theta_6,\]

where the six coordinates correspond to the six substrates of §4:
instructions, skills, memory, tools, orchestration, and governance. A
coordinate may be a line in a \texttt{CLAUDE.md} file, a Cursor rule, a
skill file, a memory entry, a retrieval configuration, a tool schema, an
agent-routing edge, an eval threshold, or a promotion rule. The space
\(\Theta_s\) is not smooth. It is discrete, symbolic, and versioned. But
it is still a parameter space: changing \(\theta_s\) changes the
behavior of the model--scaffold composite.

For a deployment patch \(D\), define patch loss as

\[L_D(\theta_s) = \mathbb{E}_{x \sim D}\left[\ell(\pi(\theta_M, \theta_s, x))\right],\]

where \(\pi(\theta_M,\theta_s,x)\) is the output or action trace
produced by the frozen frontier model under scaffold \(\theta_s\), and
\(\ell\) is an operational loss. Depending on the patch, \(\ell\) may be
a hard-failure indicator, user-correction rate, eval-suite error,
tool-call failure, retrieval miss, schema violation, reviewer
disagreement, or graded task score. The important distinction is
substrate: \(\theta_M\) is fixed during deployment; \(\theta_s\) is the
object being fit.

A failure observation does not directly yield an analytic gradient. It
yields a candidate edit. Let

\[z_t \in \mathcal{A}(\theta_s^t)\]

be an admissible scaffold edit proposed at time \(t\), drawn from a
proposal process

\[z_t \sim Q_t(z \mid H_t, \theta_s^t),\]

where \(H_t\) is the accumulated history of failures, corrections,
rejected outputs, eval regressions, reviewer comments, tool-call traces,
and prior accepted deltas, and conditioning on \(\theta_s^t\) ensures
proposals are admissible at the current scaffold. The admissible edit
set \(\mathcal{A}(\theta_s^t)\) is the local neighborhood of the current
scaffold: the rule can be rewritten, a skill added, a memory corrected,
a tool schema tightened, an orchestration edge changed, an eval case
inserted, a promotion policy modified.

Because the space is discrete, the relevant gradient analogue is not
\(\nabla_{\theta_s} L_D\). It is the \emph{finite directional
difference} in loss induced by an edit:

\[\Delta_z L_D(\theta_s) = L_D(\theta_s \oplus z) - L_D(\theta_s),\]

where \(\theta_s \oplus z\) denotes the scaffold obtained by applying
edit \(z\). A useful edit is one for which

\[\Delta_z L_D(\theta_s) < 0.\]

The gate estimates whether this inequality holds. Let
\(\widehat{\Delta}_z L_D\) be the gate's noisy estimate of the
directional loss change. In an eval-gated system,
\(\widehat{\Delta}_z L_D\) is measured on a held-out task suite. In an
LLM-judge system, it is estimated by model-mediated scoring, ideally
calibrated against human agreement. In a reviewer-gated system, it is
estimated by human credit assignment and PR review. In an RL-threshold
system, it is estimated by reward improvement. In a verifier-gated
system, it is estimated by a symbolic, learned, or surrogate verifier.

The update rule is therefore

\[\theta_s^{t+1} = \begin{cases} \theta_s^t \oplus z_t, & \text{if } G_t(z_t)=1, \\ \theta_s^t, & \text{if } G_t(z_t)=0, \end{cases}\]

where the gate \(G_t\) accepts when the estimated improvement clears the
deployment's threshold:

\[G_t(z_t)=1 \quad \Longleftrightarrow \quad \widehat{\Delta}_{z_t} L_D(\theta_s^t) + \lambda\,\Omega(z_t) \leq -\tau_t.\]

Here \(\Omega(z_t) \geq 0\) is an optional complexity or risk penalty
--- longer instructions, broader tool permissions, larger memory writes,
more invasive topology rewrites, wider promotion scope --- while
\(\lambda \geq 0\) controls regularization and \(\tau_t \geq 0\) is the
required improvement margin. Production systems often implement this
without writing the equation: an eval suite blocks a merge if regression
appears; a reviewer rejects a vague rule; a governance system prevents a
local artifact from becoming shared until evidence accumulates.

Two distinct quantities sit in the optimization analogy and the paper
keeps them separate. \textbf{Edit magnitude} --- how much a single
accepted edit changes the local scaffold under \(\oplus\) --- plays the
role of \emph{learning rate}: a one-line rule patch is a small step, a
full rubric rewrite is a large step, a tool-schema overhaul is a large
step. \textbf{Promotion radius} --- how many downstream patches the
accepted edit affects --- controls \emph{which loss is being descended
on}: USER-level acceptance descends on one patch's loss \(L_D\); PROJECT
promotion descends on an average
\(\tfrac{1}{|D_{\text{proj}}|}\sum_{D' \in D_{\text{proj}}} L_{D'}\)
across a project's patches; STACK promotion widens the averaging set;
CORE or cross-tenant promotion descends on a federated average over many
tenants' losses. Learning rate (edit magnitude) and update radius
(promotion scope) are independently controllable, and the gate must be
calibrated against both. Governance is therefore not administrative
overhead: it is the regularization system that prevents a
small-magnitude edit from being descended on the wrong loss.

The correspondence to standard optimization is therefore not decorative:

{\def\LTcaptype{none} % do not increment counter
\begin{longtable}[]{@{}
  >{\raggedright\arraybackslash}p{(\linewidth - 2\tabcolsep) * \real{0.2683}}
  >{\raggedright\arraybackslash}p{(\linewidth - 2\tabcolsep) * \real{0.7317}}@{}}
\toprule\noalign{}
\begin{minipage}[b]{\linewidth}\raggedright
Optimization concept
\end{minipage} & \begin{minipage}[b]{\linewidth}\raggedright
Scaffold analogue
\end{minipage} \\
\midrule\noalign{}
\endhead
\bottomrule\noalign{}
\endlastfoot
Parameter vector & Instructions, skills, memories, tools, routing
graphs, eval gates, governance rules \\
Forward pass & Model--scaffold composite executing on patch input \\
Loss & Failure, user correction, eval regression, retrieval miss, tool
error, schema violation \\
Finite directional difference &
\(\Delta_z L_D(\theta_s) = L_D(\theta_s \oplus z) - L_D(\theta_s)\), the
loss change induced by a candidate edit \\
Gradient estimate & Credit assignment identifying the scaffold
coordinate responsible for the failure \\
SGD step & Accepted scaffold edit \\
Aggregated update & Multiple recurring failures aggregated before update
(either by sampling --- \emph{mini-batch averaging} --- or by
sign-of-direction agreement across \(k\) proposals ---
\emph{multi-sample voting}) \\
Learning rate & Edit magnitude under \(\oplus\) (size of a single
accepted edit) \\
Update radius & Promotion scope --- which loss is being descended on
(USER, PROJECT, STACK, TENANT, CORE) \\
Validation loss & Held-out patch eval before merge \\
Regularization & Complexity penalty \(\Omega(z)\), bloat control,
cross-patch regression checks \\
Overfitting & Scaffold improves on this patch and degrades elsewhere \\
Rollback & Reverting a harmful accepted update \\
\end{longtable}
}

\textbf{Credit assignment without chain rule.} Scaffold systems do not
compute chain-rule derivatives through markdown files or tool
registries, and the correspondence above deliberately omits a
``backpropagation'' row. They do, however, perform \emph{credit
assignment} through a pipeline: if the final answer fails, the
responsible coordinate may be a retrieval configuration three layers
upstream, a missing tool schema constraint, an underspecified
instruction, a stale memory, or a bad orchestration edge. Engineers do
this manually; LLM-judges and surrogate verifiers do it
semi-automatically. The operation is \emph{attribution} --- output error
is assigned to an upstream coordinate --- rather than backpropagation.
The descent argument of §3.2 needs only that this attribution succeeds
with non-zero probability per failure (assumption (iv) of the scope
box), not that it is performed by chain rule.

Several concrete cases make the mechanism visible.

In a RAG system, suppose failures recur because retrieved context omits
the decisive clause. The observed loss is not merely ``the model
hallucinated.'' Credit assignment traces the error to the retrieval
coordinate: chunking, query rewrite, metadata filter, embedding model,
reranker, citation policy, or context-packing rule. The candidate edit
might tighten chunk boundaries, add a metadata constraint, change the
reranker, or insert a regression eval containing the missed clause. The
gate tests whether retrieval recall and answer accuracy improve on the
patch eval without degrading adjacent cases. If accepted, the scaffold
has moved downhill on the local loss surface.

In a pitch-review workflow, suppose the system repeatedly overrates
decks that claim a ``huge market'' without numeric evidence. The failure
cluster identifies an under-specified rubric coordinate. The candidate
edit changes ``evaluate TAM size'' into ``require a numeric market
estimate and source; if absent, cap TAM score at 4/10.'' The gate runs
the revised rubric against held-out decks with human scores. If
agreement improves without damaging unrelated criteria, the edit is
committed. The rubric is not commentary. It is a scaffold parameter
updated in response to loss.

In a tool-using agent, suppose calls repeatedly fail because optional
parameters are underspecified. The responsible coordinate may be the
tool schema, the call policy, or the clarification rule. The candidate
edit adds required fields, constrained decoding, or a pre-call
uncertainty check. The gate estimates whether tool-call failure
decreases without increasing refusal, latency, or unnecessary
clarification. If accepted, the scaffold has fit a recurring local
failure mode.

In a memory system, suppose a stale customer preference is repeatedly
retrieved and causes bad decisions. The failure is not a weight error.
The responsible coordinate is a memory entry plus its provenance,
priority, and expiry policy. The candidate edit corrects, expires, or
version-tags the memory. The gate checks whether the correction improves
future behavior without erasing useful history. Again: loss, credit
assignment, candidate update, validation, commit.

Under bounded loss and a few additional regularity conditions, this is
descent in the formal sense. The critical bookkeeping is that the gate
sees an estimator \(\widehat{\Delta}_{z_t} L_D\) of the directional loss
change, not the true \(\Delta_{z_t} L_D\); acceptance is conditional on
the estimator, and the descent argument must therefore distinguish the
two.

\begin{quote}
\textbf{Assumption A1 (positive descent bias on accepted edits).} There
exists \(\gamma_t > 0\) such that

\[\mathbb{E}\!\left[\widehat{\Delta}_{z_t} L_D(\theta_s^t) \mid G_t(z_t)=1,\theta_s^t\right] \leq -\gamma_t - \tau_t.\]

The gate's accept-threshold \(\tau_t\) ensures accepted edits clear the
margin under the estimator.

\textbf{Assumption A2 (bounded estimator bias on accepted edits).} There
exists \(\varepsilon_t \geq 0\) such that

\[\big|\mathbb{E}\!\left[\Delta_{z_t} L_D - \widehat{\Delta}_{z_t} L_D \mid G_t(z_t)=1,\theta_s^t\right]\big| \leq \varepsilon_t.\]
\end{quote}

Combining A1 and A2 with the update rule and writing
\(p_t = \Pr(G_t(z_t)=1 \mid \theta_s^t)\),

\[\mathbb{E}\!\left[L_D(\theta_s^{t+1}) \mid \theta_s^t\right] \leq L_D(\theta_s^t) - p_t(\gamma_t + \tau_t) + p_t \varepsilon_t.\]

This is descent under a noisy estimator: the per-step expected
improvement is the gate's margin minus its estimator bias, weighted by
the acceptance probability.

\begin{quote}
\textbf{Proposition (convergence of loss, informal).} Assume A1, A2,
bounded loss \(0 \leq L_D(\theta_s) \leq B\), and additionally:

\begin{itemize}
\tightlist
\item
  \textbf{A3 (proposal coverage).} At every non-\(G\)-stable scaffold
  there exists \(\delta > 0\) such that the proposal distribution
  \(Q_t\) places probability at least \(q > 0\) on an admissible edit
  with \(\mathbb{E}\!\left[\Delta_z L_D\right] \leq -\delta\).
\item
  \textbf{A4 (effective compactness).} The reachable scaffold orbit
  \(\{\theta_s^t\}\) takes values in a finite or precompact subset of
  \(\Theta_s\).
\item
  \textbf{A5 (persistent acceptance with summable bias).}
  \(\sum_t p_t(\gamma_t + \tau_t) = \infty\) and
  \(\sum_t p_t \varepsilon_t < \infty\).
\end{itemize}

Then \(L_D(\theta_s^t)\) is a non-negative supermartingale up to
summable noise; \(L_D(\theta_s^t) \to L_\infty\) almost surely for some
random \(L_\infty \geq 0\); and the limit set of \(\{\theta_s^t\}\) is
\textbf{\(G\)-stable}: no admissible edit clears the gate with negative
expected directional loss in the limit.
\end{quote}

The proposition claims convergence of the \emph{loss} and stability of
the \emph{gate decision}, not convergence to a global or even local
optimum of \(L_D\). A globally suboptimal \(G\)-stable scaffold is
consistent with the result --- for example, when the proposal
distribution \(Q_t\) never proposes the improving edit that would unlock
further descent, or when the estimator \(\widehat{\Delta}\) is biased in
a way that consistently rejects an improving direction. This is the
ordinary target of local stochastic optimization under noisy estimators
rather than a guarantee of optimality.

In non-stationary patches, \(D = D_t\), the same mechanism tracks a
moving optimum rather than converging once. This is why recency
weighting, replay, pruning, expiry, and rollback are not optional
hygiene. They are the scaffold-level analogues of continual-learning
machinery. Without them, accumulated edits create scaffold bloat, stale
rules, memory pollution, and patch-level plasticity loss.

This formalization also clarifies the role of overfitting. In weight
training, overfitting is usually a defect: the model becomes too
specialized to the training distribution and loses transfer. In scaffold
adaptation, local overfitting is partly the point. Frontier weights are
trained to preserve cross-patch generality; they must smooth over local
conventions, rare workflow-specific failures, contradictory tenant
preferences, and hidden evaluator expectations. A deployment patch needs
the opposite: selective specialization to local residuals. The scaffold
is therefore an intentionally overfittable layer, but one whose
overfitting is scoped, inspectable, versioned, gated, and reversible. A
scaffold that improves one repository, hospital workflow, legal
diligence process, or support desk may degrade performance elsewhere.
That is not a contradiction. It is patch adaptation. The mistake is
promoting that edit beyond its evidence radius.

Loop 1 and Loop 2 are now mathematically interpretable, and they map
cleanly onto the \emph{edit-magnitude} / \emph{update-radius} split of
the analogy table. Loop 1 proposes and validates local descent steps and
operates on the per-patch loss \(L_D\). Loop 2 expands the update radius
--- it selects, merges, and promotes accepted local edits so that they
apply to a larger averaging set of patch losses (a project, a stack, a
tenant, or a cross-tenant pool). USER-level acceptance descends on a
single \(L_D\). PROJECT-level promotion descends on a project-averaged
loss. STACK-level promotion widens the average further. CORE or
cross-tenant promotion descends on a federated average. The two-loop
design space is therefore not just a taxonomy of engineering patterns:
it is a taxonomy of \emph{which loss each surveyed system is descending
on}, gated by which estimator.

The artifact-layer descent framing is mechanism when five conditions
hold: measurable patch loss, residual capture, failure-to-coordinate
credit assignment, candidate delta synthesis, and a gate that estimates
directional loss change before persistence or promotion. It is merely
metaphor when those conditions are absent. This distinction explains the
survey split. Research full-auto systems automate delta synthesis and
validation but often lack governance. Production governance systems
validate and promote but often lack failure-triggered candidate
generation. Hybrid systems close more of the loop. The empty enterprise
corner is the case where all pieces must hold across tenants.

The conclusion is simple: the frontier model is trained to avoid
overfitting across patches. The scaffold exists to overfit safely inside
one.

\subsection{4. Six scaffold substrates}\label{six-scaffold-substrates}

A patch-plastic scaffold composes six substrates. They are the
coordinates of \(S_t\).

{\def\LTcaptype{none} % do not increment counter
\begin{longtable}[]{@{}
  >{\raggedright\arraybackslash}p{(\linewidth - 4\tabcolsep) * \real{0.3333}}
  >{\raggedright\arraybackslash}p{(\linewidth - 4\tabcolsep) * \real{0.3333}}
  >{\raggedright\arraybackslash}p{(\linewidth - 4\tabcolsep) * \real{0.3333}}@{}}
\toprule\noalign{}
\begin{minipage}[b]{\linewidth}\raggedright
Substrate
\end{minipage} & \begin{minipage}[b]{\linewidth}\raggedright
What persists
\end{minipage} & \begin{minipage}[b]{\linewidth}\raggedright
Score-5 mark
\end{minipage} \\
\midrule\noalign{}
\endhead
\bottomrule\noalign{}
\endlastfoot
\textbf{S1 --- Instructions} & Project/team/org rules attached to a
workspace & Versioned, two-loop, Auto-gated promotion \\
\textbf{S2 --- Skills} & Procedural artifacts (code/text) loaded on
demand & Skill bank self-modifies on failure, unit-test gated,
versioned, promoted \\
\textbf{S3 --- Memory} & Cross-session experience (episodic / semantic /
procedural) & Provenance + versioning + correction pathways \\
\textbf{S4 --- Tools} & Vetted action surface + context schemas &
Domain-curated bundle is the product; eval-gated tool versions \\
\textbf{S5 --- Orchestration} & Multi-role graph + routing &
Population/policy updates from outer-loop signal \\
\textbf{S6 --- Governance} & Versioning, promotion, audit, rollback &
dev→staging→prod with eval thresholds + audit log + rollback \\
\end{longtable}
}

The 0--5 rubric is uniform across substrates: \textbf{0} ephemeral;
\textbf{1} persistent local; \textbf{2} persistent + tool attachment, no
feedback; \textbf{3} multi-scope or feedback but no automated promotion;
\textbf{4} automated feedback updates the scaffold (one loop closed);
\textbf{5} two-loop versioned promotion plus governance (both loops
closed).

Each substrate maps to a different Paper-2 residual class. S1 and S4
carry capability provisioning. S2 and S3 accumulate the patch-specific
failure-mode catalogue. S5 lowers per-step hard-fraction by
specialisation. S6 is the gate substrate --- without it, none of the
other five can be updated safely across deployments. \emph{The
substrates are not a taxonomy of nice-to-haves. They are the layers of
the scaffold parameter vector, separately addressable by Loop 1 and Loop
2.} Score distributions and per-substrate exemplars are summarised in
§6.1; the maturity-rubric details for practitioner adoption are deferred
to follow-on work. We adopt this six-coordinate decomposition for the
survey because each coordinate is independently editable in production
systems (a Cursor rule can change without touching a tool schema; an
eval threshold can change without touching a memory store); alternative
partitions are possible (e.g., separating evals from governance, or
merging instructions and skills) and the descent argument of §3.2 is
invariant to the partition so long as the independent-editability
property holds.

\subsection{5. Survey method and evidence
tiers}\label{survey-method-and-evidence-tiers}

The survey covers approximately 130 unique LLM systems published or
productionised between January 2023 and May 2026, split into two
sets.\footnote{The master CSV
  (\texttt{supplementary/p3\_master\_scores.csv}) records each system as
  one row per scored substrate, so the row count exceeds the
  unique-system count. Cited figures throughout §§5--8 refer to unique
  systems; Appendix B reconciles the per-substrate row count, the
  unique-system count, and the patch-plastic core/contrast split.} The
\textbf{core corpus} of approximately 90 candidate scaffold-learning
systems satisfies the patch-plastic discriminator (score ≥ 3 on at least
one substrate, indicating skills, memory, or multi-scope persistence
beyond ephemeral prompts). The \textbf{contrast corpus} of approximately
38 systems is commonly described in the public discourse as agentic or
self-improving but fails at least one discriminator --- typically Loop 2
absent, no persistent artifact distinct from optimiser state, or
in-context-only adaptation. Headline statistics quote the core corpus;
the contrast set is preserved as asterisked entries in the matrices for
definitional clarity.

\textbf{Inclusion.} Any system that (a) wraps a foundation model with at
least one persistent artifact updated from deployment signal, or (b) is
named in the public discourse as a ``self-improving,'' ``scaffolded,''
or ``agentic'' system. The contrast set (Self-Refine, ReAct, plain Mem0,
DSPy/TextGrad, etc.) is included precisely so the discriminator has
something to mark against.

\textbf{Exclusion.} Systems whose only update mechanism is weight-level
fine-tuning or RLHF --- these are model-side, not scaffold-side, and
belong to the self-evolution literature this paper differentiates from.
Systems with no public evidence of deployment (white papers without
code, demos, or production claim).

\textbf{Evidence tiers.} Each system carries one or more of: (T1)
peer-reviewed publication, (T2) arXiv preprint, (T3) production claim
with metrics, (T4) open-source artifact (GitHub, registry), (T5)
industry-blog or interview-level documentation. Systems with T1 or T4
evidence carry highest weight in the cluster analysis (§7); T3 systems
are admissible where the claim is concrete (specific benchmark, named
integration target); T5 systems are admissible only when corroborated by
T1--T4 elsewhere. Tier flags are propagated to cluster claims through
the per-row evidence column in Appendix A; T3 production claims are not
pooled with T1/T2 evidence when computing cluster-level statistics.
\textbf{Survey window and dating.} The corpus includes arXiv preprints
published through May 2026; preprint-only systems carry T2 evidence and
are flagged in Appendix A. A fraction of T2 work is recent enough that
ablations may yet be revised; cluster-level claims are robust to
dropping any single T2 system.

\textbf{Scoring.} Each system received a substrate-by-substrate 0--5
score against the rubric of §4 and an integrated patch-plasticity (PP)
score in \{0,1,2,3,3.5,4,4.5,5\}. Half-scores were used sparingly for
systems where one criterion was clearly met and another partially.
Ambiguous cases were resolved conservatively (downward). Scoring is
single-rater; the supplementary CSV exposes every per-system score so
reviewers can audit ambiguous classifications independently against the
per-substrate evidence harvests
(\texttt{supplementary/p3\_harvest\_s1\_instructions.md} through
\texttt{p3\_harvest\_s6\_governance.md}).

\textbf{Composite two-loop criterion.} A separate, stricter audit
applied to 22 candidate systems. A system qualifies as composite
two-loop if and only if it satisfies both: Loop 1 modifies a persisted
artifact from a session signal \emph{and} Loop 2 has an explicit
cross-context promotion mechanism with versioning or lineage. Thirteen
research systems passed the strict audit, with two additional industry
exemplars (Sierra Agent OS 2.0, Cognition Devin) evaluated using vendor
documentation as primary evidence and explicitly flagged in the matrix
of §8; nine systems failed (MAE, MetaGen, MetaReflection, CLIN ---
Loop-2 absent; DSPy, TextGrad --- no persistent artifact distinct from
optimiser state; Mem0, Zep, ChatGPT Memory --- no Loop-2 promotion
path).

\textbf{Sensitivity analysis.} Under relaxation (a) --- counting any of
\emph{versioning OR lineage OR rollback} as the Loop-2 governance
criterion --- the composite-two-loop count rises from thirteen research
systems to approximately seventeen, by re-admitting MAE, MetaGen,
MetaReflection, and CLIN that were excluded for lacking versioned
promotion lineage. Under relaxation (b) --- admitting session signal
that updates optimizer state without persisting an artifact distinct
from optimiser state --- the count rises to approximately fifteen, by
re-admitting DSPy and TextGrad. The qualitative findings (Auto×Auto
well-populated, {[}Human × Human{]} and {[}Human × Auto{]} empty in the
surveyed corpus, governance-first cluster in {[}None × Auto{]},
enterprise full-auto cell empty) survive both relaxations. The CSV
bundled in \texttt{supplementary/} permits readers to re-run either
relaxation.

\textbf{Audit trail.} The full per-system spreadsheet
(\texttt{supplementary/p3\_master\_scores.csv}), per-substrate evidence
harvests (\texttt{supplementary/p3\_harvest\_s1\_instructions.md}
through \texttt{p3\_harvest\_s6\_governance.md}), integrated audit
(\texttt{supplementary/p3\_harvest\_integrated.md}), composite-two-loop
audit (\texttt{supplementary/p3\_harvest\_composite\_two\_loop.md}),
failure-mode harvest
(\texttt{supplementary/p3\_harvest\_failure\_modes.md}), and prior-art
check (\texttt{supplementary/p3\_prior\_art\_check.md}) ship inside the
arXiv source bundle. Reviewers can re-run the analysis from primary
sources.

\subsection{6. Survey findings: the scaffold is unevenly
mature}\label{survey-findings-the-scaffold-is-unevenly-mature}

The patch-plastic core corpus does not mature uniformly. Tools and
integrated systems are the most mature substrates; skills and
orchestration lag; instructions are ubiquitous but shallow; memory is
the most paradoxical of the six. This section presents the headline
distributions before the cluster interpretation of §7. The full
per-system spreadsheet is \texttt{supplementary/p3\_master\_scores.csv};
Appendix A lists representative rows.

\subsubsection{6.1 Substrate maturity is
uneven}\label{substrate-maturity-is-uneven}

Aggregating the survey by substrate gives this distribution:

{\def\LTcaptype{none} % do not increment counter
\begin{longtable}[]{@{}
  >{\raggedright\arraybackslash}p{(\linewidth - 8\tabcolsep) * \real{0.1400}}
  >{\raggedleft\arraybackslash}p{(\linewidth - 8\tabcolsep) * \real{0.0400}}
  >{\raggedleft\arraybackslash}p{(\linewidth - 8\tabcolsep) * \real{0.0500}}
  >{\raggedright\arraybackslash}p{(\linewidth - 8\tabcolsep) * \real{0.2200}}
  >{\raggedright\arraybackslash}p{(\linewidth - 8\tabcolsep) * \real{0.5500}}@{}}
\toprule\noalign{}
\begin{minipage}[b]{\linewidth}\raggedright
Substrate
\end{minipage} & \begin{minipage}[b]{\linewidth}\raggedleft
\(n\)
\end{minipage} & \begin{minipage}[b]{\linewidth}\raggedleft
Mean
\end{minipage} & \begin{minipage}[b]{\linewidth}\raggedright
Score-5 systems
\end{minipage} & \begin{minipage}[b]{\linewidth}\raggedright
Survey finding
\end{minipage} \\
\midrule\noalign{}
\endhead
\bottomrule\noalign{}
\endlastfoot
\textbf{S1 --- Instructions} & 12 & 2.83 & \emph{(none)} & Rules persist
across IDEs and agents (Claude Code's CLAUDE.md, Cursor Rules,
AGENTS.md, Windsurf, Continue.dev) but no surveyed instruction system
closes both an automated inner loop and a governed two-loop promotion.
Ceiling held at 4 by Claude Code and Replit \\
\textbf{S2 --- Skills} & 16 & 2.62 & SAGE, COSPLAY & Research frontier
is active (failure-triggered skill update); production governance is
weaker. Anthropic's skills repository is one-pool with PR review but no
automated eval gate \\
\textbf{S3 --- Memory} & 19 & 3.16 & Cuadros et al.~governed
collaborative memory & Cross-session memory is widespread;
\emph{correctable, versioned} memory is rare. Production memory systems
(ChatGPT Memory, Claude Memory, Mem0, Zep, MemGPT) treat memory as
append-only \\
\textbf{S4 --- Tools} & 19 & 3.63 & MCP, Harvey, Hippocratic AI & Most
mature production substrate. Tool bundles are the commercial unit; MCP
is the cross-vendor standard with broad public-server ecosystem and
substantial enterprise adoption \\
\textbf{S5 --- Orchestration} & 19 & 2.68 & \emph{(none)} & Multi-agent
topologies exist (LangGraph, AutoGen, CrewAI) but topology
\emph{evolution} is rare; MAE comes closest with RL co-evolution of
Proposer/Solver/Judge population but lacks Loop-2 cross-context
promotion and so does not satisfy the score-5 rubric; production rarely
exposes governed promotion paths for emerging agent graphs \\
\textbf{S6 --- Governance} & 21 & 3.05 & Braintrust, Vellum, LangSmith
Hub, AGENTS.md/AAIF & High production maturity for prompts-as-code ---
immutable commits, eval-gated PR merge, sub-five-minute rollback ---
typically without inner-loop adaptation \\
\textbf{INTEGRATED} & 34 & 3.74 & NanoResearch, AutoAgent, SkillRL & The
high-water mark when present, but the cluster is dominated by research
systems lacking enterprise governance \\
\end{longtable}
}

The unevenness is the survey's central empirical pattern. The field has
independently converged on the six substrates, but no system matures all
six at once. Research systems learn fast and lack enterprise governance;
production systems govern well and learn slowly; open standards solve
artifact portability but not adaptation; vertical-bundle vendors solve
patch specificity but rarely expose cross-tenant promotion. The missing
architecture is not ``agents with tools'' or ``memory with evals.'' It
is \emph{governed scaffold adaptation} --- fast local learning paired
with safe cross-context promotion.

\subsubsection{6.2 Representative systems by
cluster}\label{representative-systems-by-cluster}

Four clusters recur across the corpus. The grouping below maps each
system to its gate types and to the structural lesson it carries. The
matrix in §8 places the same systems on the 3×3 design space.

{\def\LTcaptype{none} % do not increment counter
\begin{longtable}[]{@{}
  >{\raggedright\arraybackslash}p{(\linewidth - 6\tabcolsep) * \real{0.1400}}
  >{\raggedright\arraybackslash}p{(\linewidth - 6\tabcolsep) * \real{0.3000}}
  >{\raggedright\arraybackslash}p{(\linewidth - 6\tabcolsep) * \real{0.1400}}
  >{\raggedright\arraybackslash}p{(\linewidth - 6\tabcolsep) * \real{0.4200}}@{}}
\toprule\noalign{}
\begin{minipage}[b]{\linewidth}\raggedright
Cluster
\end{minipage} & \begin{minipage}[b]{\linewidth}\raggedright
Representative systems (year, PP score)
\end{minipage} & \begin{minipage}[b]{\linewidth}\raggedright
Loop 1 / Loop 2
\end{minipage} & \begin{minipage}[b]{\linewidth}\raggedright
What the cluster proves
\end{minipage} \\
\midrule\noalign{}
\endhead
\bottomrule\noalign{}
\endlastfoot
Research full-auto & NanoResearch (2026, 5), SkillRL (2026, 5),
AlphaEvolve (2025, 4), DGM (2025, 4), ADAS (2024, 4), Voyager (2023, 4)
& Auto / Auto & Both loops can be fully automated under algorithmic
gates; benchmark gains compound; enterprise governance is uniformly
absent \\
Production hybrid & SkillForge (2026, 4), CASCADE (2025, 4.5), Sierra OS
2.0 (2025, 3)\footnote{Sierra Agent OS 2.0 also appears in the
  \emph{Open standards / vertical-bundle} discussion in §7 because its
  commercial product wraps a vertical workflow around the standardised
  agent-OS surface; we place it primarily in \emph{Production hybrid} on
  the basis of its Loop-1/Loop-2 gate types.}, Cognition Devin (2024, 3)
& Auto / Human & Algorithmic Loop-1 gate plus reviewer-judgment Loop-2
gate is the production-viable cell; SkillForge is the architectural
exemplar \\
Governance-first & Braintrust, Vellum, LangSmith Hub (all 2024, S6 = 5)
& None / Auto & Eval-gated promotion exists; candidate generation is
manual (engineers iterate prompts); no patch-plastic inner loop \\
Open standards & AGENTS.md (2025), MCP (2024), Claude Skills (2025),
Cursor Rules (2024) & Mixed / Mixed & Vendor-cross convergence on
artifact format and tool attachment; tens of thousands of AGENTS.md
repositories; broad MCP server ecosystem with substantial enterprise
adoption \\
\end{longtable}
}

\subsubsection{6.3 Surprising absences}\label{surprising-absences}

What is missing from the survey is as instructive as what is present.
Five absences stand out across the patch-plastic core corpus:

\begin{enumerate}
\def\labelenumi{\arabic{enumi}.}
\tightlist
\item
  \textbf{Failure-triggered memory is rare in production.} Deployed
  memory systems prefer recording preferences and successful facts; only
  Reflexion, the Voyager skill library, Gemini Code Assist's
  PR-rejection memory, and the classical SOAR chunking pattern
  explicitly record residuals. The most obvious Loop-1 signal is the one
  production memory systems mostly ignore.
\item
  \textbf{Cross-customer skill promotion is absent.} AGENTS.md is
  cross-vendor; Anthropic's skills repository is one-pool. No surveyed
  system has a vetted cross-organisational marketplace with eval-gated
  promotion.
\item
  \textbf{Memory versioning is mostly unsolved.} Only the Cuadros et
  al.~(2026) governed collaborative memory framework provides explicit
  provenance, versioning, and correction pathways. Production default is
  append-only with implicit retrieval-based forgetting --- at odds with
  the patch-plastic discriminator's correction requirement.
\item
  \textbf{Scaffold-level curriculum is absent.} No surveyed system
  selects \emph{which} sessions inform the gradient. Every session
  contributes equally; weighting by recency, severity, or
  representativeness has not been published.
\item
  \textbf{Overfitting detection at Loop-2 promotion is absent.} No
  surveyed system runs a held-out eval set per project to catch
  scaffolds that over-fit one deployment before they propagate to
  another.
\end{enumerate}

Each absence is a Paper-4 target. The conclusion echoes them; the survey
makes them visible.

\subsection{7. Convergent patterns}\label{convergent-patterns}

The corpus, scored against the rubric of §4, produces four archetypal
clusters that recur across both research and industry. We present the
clusters here rather than walking every system; the full table is in
\texttt{supplementary/p3\_master\_scores.csv}.

\textbf{Research full-auto.} Systems where both loops fire from
algorithmic criteria and no human gates the routine update.
\textbf{NanoResearch} (Xu et al., 2026) co-evolves Skills, Memory, and
Policy with a semantic-merge orchestrator; Innovation 4.96 → 5.65 and
Compliance 6.66 → 8.96 across three rounds. \textbf{SkillRL}
(arXiv:2602.08234) uses RL reward distillation with a hard threshold at
task-category accuracy below 0.4 and reports +15.3\% over strong
baselines. \textbf{AlphaEvolve} (DeepMind, 2025) runs a Gemini-driven
ensemble with a cascade evaluator and a MAP-Elites archive, deployed at
Google to gain +0.7\% Borg compute worldwide, +23\% on the Gemini
training kernel, +32.5\% on FlashAttention, and the first general-field,
recursively applicable improvement over Strassen's 1969 bound on 4×4
complex matrix multiplication (48 scalar multiplications). The honest
framing matters: AlphaTensor (Fawzi et al., 2022) had already found a
47-multiplication algorithm over GF(2), and Winograd reported a
48-multiplication method for commutative rings that was not recursively
applicable; AlphaEvolve's contribution is the first improvement that
holds in general fields and admits recursive composition. \textbf{DGM}
(Zhang et al., 2025) modifies its own source code with archive-based
parent selection; SWE-bench 20\% → 50\%. The cluster maximises velocity.
Enterprise governance is uniformly absent.

\textbf{Production hybrid.} Systems where Loop 1 is Auto-gated (an
algorithmic criterion) and Loop 2 is Human-gated (reviewer approval
before propagation). \textbf{SkillForge} (Alibaba Cloud,
arXiv:2604.08618) drives auto-generated skill diffs through an LLM-judge
with \textgreater90\% human agreement (the Loop-1 gate); human support
engineers then review what the LLM-judge passes (the Loop-2 gate). VFS
version commits provide full lineage. \textbf{CASCADE} (Huang et al.,
2025) reports the largest within-benchmark ablation gain in the surveyed
corpus --- 93.3\% vs.~35.4\% on SciSkillBench (+57.9 pp) --- driven by
memory consolidation gated by human-agent collaboration. The headline
gain is not directly comparable to AlphaEvolve's production +0.7\% Borg
or DGM's +30 pp on the broader SWE-bench: it is measured on a single
benchmark designed to reward the very mechanism CASCADE adds, and should
be read as evidence of mechanism-on-benchmark fit rather than as a
cross-cluster effect size. \textbf{Sierra Agent OS 2.0} routes AI-driven
Insights through human-authored Expert Answers in GitHub-style
Workspaces. \textbf{Cognition Devin} ships dynamic in-session tool
creation with a 67\% PR-merge rate (the merged PRs are the Loop-2
promotion). This is the production-viable cell.

\textbf{Governance-first.} Systems where Loop 2 is Auto-gated by an eval
criterion but Loop 1 is essentially absent --- there is no automatic
candidate-generation pipeline. \textbf{Braintrust}, \textbf{Vellum},
\textbf{LangSmith Hub}, \textbf{W\&B Weave}, \textbf{Agenta},
\textbf{LangFuse}. The pattern is \emph{prompts-as-code}: immutable
commits, semantic-version tag pointers, PR-blocked merges on eval
regression, sub-five-minute rollback. Engineers iterate the artifacts
manually; the eval pipeline gates the promotion. We frame this as
\textbf{MLOps for scaffolds}, distinct from MLOps for weights.

\textbf{Open standards.} Systems where the artifact format and the
layering convention are open standards across vendors.
\textbf{AGENTS.md} (Linux Foundation, December 2025) is now in over
60,000 GitHub repositories; Lulla et al.~(2026) report −28.64\%
wall-clock and −16.58\% tokens with no quality loss after adding an
AGENTS.md scaffold. \textbf{MCP} (Anthropic, 2024) has become the
cross-vendor standard with thousands of public servers and substantial
enterprise adoption. \textbf{Claude Agent Skills}, \textbf{Cursor
Rules}, \textbf{GitHub Copilot custom instructions}, \textbf{Windsurf
rules}, \textbf{Continue.dev rules}, \textbf{Zed AI rules},
\textbf{Replit \texttt{replit.md}} all converge on the same pattern: a
git-tracked markdown artifact attached to a workspace, with an MCP-style
tool layer for actions. Convergence across vendors is the strongest
evidence the architecture is real rather than a single vendor's local
optimum.

The four clusters do not partition the corpus --- many systems sit
between two clusters --- but they capture the architectural variation.
Vertical-bundle vendors (Harvey for legal, Hippocratic AI for
healthcare, Sierra for customer experience) sit between \emph{open
standards} and \emph{production hybrid}: the bundle is the commercial
unit; the underlying model is mostly the same one a competitor would
use.

One result from the orchestration literature bounds how strongly any of
these clusters can be read as evidence for multi-agent architectures
specifically. Tran \& Kiela (2026) report that single-agent LLMs match
or beat multi-agent systems on multi-hop reasoning when thinking-token
budgets are matched; the implication is that observed maturity gaps on
S5 (orchestration) may reflect the rarity of governed topology
\emph{evolution} and the conflation of agent-count with compute-budget
rather than a structural deficit of multi-agent designs. The
orchestration cluster is therefore better read as evidence that
production has converged on flat agent topologies and on tool-mediated
rather than role-mediated decomposition, not as evidence that
orchestration is the binding maturity constraint.

\subsection{8. The two-loop design
space}\label{the-two-loop-design-space}

The composite-two-loop systems sit in a 3×3 matrix indexed by gate type
on each loop. Thirteen research systems satisfy the strict criterion,
with two production exemplars (Sierra Agent OS 2.0, Cognition Devin)
included on the basis of vendor documentation as primary evidence and
explicitly flagged.

{\def\LTcaptype{none} % do not increment counter
\begin{longtable}[]{@{}
  >{\raggedright\arraybackslash}p{(\linewidth - 6\tabcolsep) * \real{0.2500}}
  >{\raggedright\arraybackslash}p{(\linewidth - 6\tabcolsep) * \real{0.2500}}
  >{\raggedright\arraybackslash}p{(\linewidth - 6\tabcolsep) * \real{0.2500}}
  >{\raggedright\arraybackslash}p{(\linewidth - 6\tabcolsep) * \real{0.2500}}@{}}
\toprule\noalign{}
\begin{minipage}[b]{\linewidth}\raggedright
\end{minipage} & \begin{minipage}[b]{\linewidth}\raggedright
\textbf{Loop-2 Auto-gated}
\end{minipage} & \begin{minipage}[b]{\linewidth}\raggedright
\textbf{Loop-2 Human-gated}
\end{minipage} & \begin{minipage}[b]{\linewidth}\raggedright
\textbf{Loop-2 None}
\end{minipage} \\
\midrule\noalign{}
\endhead
\bottomrule\noalign{}
\endlastfoot
\textbf{Loop-1 Auto-gated} & NanoResearch, SkillRL, AlphaEvolve, ADAS,
DGM\footnote{DGM operates in both Auto×Auto (primary archive-based
  evolution) and Auto×Human (sandboxed code-modification review) modes
  depending on deployment configuration; we list it once in the dominant
  cell.}, Voyager, SAGE, EvolveR, CoEvoSkills, AutoAgent, AutoManual &
SkillForge, CASCADE, Sierra OS 2.0†, Cognition Devin† & MAE*, MetaGen*,
MetaReflection*, CLIN* \\
\textbf{Loop-1 Human-gated} & \textbf{{[}empty by engineering logic{]}}
& \textbf{{[}empty in surveyed corpus{]}} & \emph{(none)} \\
\textbf{Loop-1 None} & ExpeL (weak), Braintrust*, LangSmith Hub*,
Vellum* & Anthropic skills repo*, DSPy*, TextGrad* & Self-Refine*,
ReAct*, Mem0* \\
\end{longtable}
}

\emph{Asterisks fail at least one composite-two-loop criterion; included
for contrast.} † Industry exemplar evaluated using vendor documentation
as primary evidence.

The matrix has three populated cells of interest. \textbf{{[}Auto ×
Auto{]}} is most populated (eleven systems). Every member has a
machine-evaluated decision criterion on both loops. \textbf{{[}Auto ×
Human{]}} is the production-viable cell --- Loop 1 fires at machine
speed under an algorithmic criterion, Loop 2 requires human approval
before propagation. \textbf{{[}None × Auto{]}} holds the
governance-first prompts-as-code systems. Three cells are empty for
distinct reasons. \textbf{{[}Human × Auto{]}} is empty by engineering
logic; if humans author candidates manually, automated promotion offers
little leverage. A speculative occupant would be ``expert curation +
auto-distribution'' (expert-authored skills automatically gated through
a held-out eval before shipping), but no published example exists.
\textbf{{[}Human × Human{]}} is empty in the surveyed corpus; a
candidate occupant would be a reviewer-authored-and-reviewer-promoted
artifact pipeline targeting solo and small-team workflows where the
marginal cost of a bad auto-merge exceeds the marginal cost of human
latency --- no public example was found during the May 2026 survey
window. \textbf{{[}Human × None{]}} is structurally improbable: a system
that requires human authorship on Loop 1 but has no Loop 2 at all
collapses to ``engineer-edits-a-config'' and is not patch-plastic in the
sense of §3.

The \textbf{{[}None × Auto-gated{]}} cell holds the governance-first
prompts-as-code systems (Braintrust, LangSmith Hub, Vellum). They have a
real Loop 2 --- eval-gated CI blocks promotion on regression --- but no
Loop 1 in the patch-plastic sense, because the candidate artifacts are
authored by engineers at their desks rather than triggered by session
feedback. The asterisks mark this: governance without inner-loop
adaptation. The cell is informative precisely because it shows what
survives when only the outer loop is automated.

Two cross-cutting findings sharpen the picture. \textbf{Noise reduction
via batching is largely absent in the surveyed auto-gated systems}:
every auto-gated system in {[}Auto × Auto{]} and {[}Auto × Human{]}
fires at \(n = 1\) per generation step, absorbing estimator noise
statistically via the per-edit gate rather than aggregating evidence
across multiple candidate edits. Multi-sample agreement filters (accept
only when \(k\) independent dreams or trajectories converge on the same
proposal) are a tractable design dimension that no surveyed system has
formalised; the noise-reduction behaviour of such a filter is
qualitatively different from mini-batch averaging --- false-accept
probability shrinks roughly as \(p^k\) rather than variance shrinking as
\(1/k\) --- and the two should not be conflated. \textbf{Multi-layer
hierarchies are under-explored}: only AlphaEvolve demonstrates explicit
depth in the auto-gated cluster, via the cascade evaluator and
MAP-Elites archive. The other twelve composite-two-loop systems treat
their artifact pool as flat. Adding a STACK-style intermediate layer to
NanoResearch or a multi-archetype scope to Anthropic's skills repository
is a tractable engineering direction.

\subsection{9. The empty corner}\label{the-empty-corner}

Four properties define what we call \textbf{enterprise full-auto}:

\begin{itemize}
\tightlist
\item
  \begin{enumerate}
  \def\labelenumi{(\alph{enumi})}
  \tightlist
  \item
    \textbf{Auto-gated Loop 1} --- the commit criterion is an
    algorithmic threshold.
  \end{enumerate}
\item
  \begin{enumerate}
  \def\labelenumi{(\alph{enumi})}
  \setcounter{enumi}{1}
  \tightlist
  \item
    \textbf{Multi-tenant cross-org promotion} --- Loop 2 routes
    artifacts between organisations, not just within one.
  \end{enumerate}
\item
  \begin{enumerate}
  \def\labelenumi{(\alph{enumi})}
  \setcounter{enumi}{2}
  \tightlist
  \item
    \textbf{Versioned lineage with RBAC} --- promoted artifacts carry
    semantic versions, audit trails, and access-control policies.
  \end{enumerate}
\item
  \begin{enumerate}
  \def\labelenumi{(\alph{enumi})}
  \setcounter{enumi}{3}
  \tightlist
  \item
    \textbf{Rollback} --- a bad promotion can be reverted without
    redeploy.
  \end{enumerate}
\end{itemize}

No surveyed system has all four. Closest approaches partition the
requirements: \textbf{AlphaEvolve} has (a) and partial (b) within a
single organisation; \textbf{NanoResearch} has (a) at the research
level; \textbf{SkillForge} has (a) and (b) within one enterprise,
partial (c) via VFS commits, undocumented (d); \textbf{Sierra Agent OS
2.0} has (b) and partial (a); \textbf{Vellum / Braintrust / LangSmith
Hub} have (c) and (d) explicitly but no (a) or (b) in the strict sense.
No surveyed system combines (b), (c), and (d) under any gate type ---
the cross-organisational governance triad is itself unfilled,
independently of whether Loop 1 is auto-gated.

The empty corner is concrete in ML terms: a \emph{DP-FedAvg-style
cross-tenant aggregator with an automated eval gate at the scaffold
layer}. Differential-privacy federated averaging (DP-FedAvg) is the
closest analogue in the federated-learning literature --- gradient
updates aggregated across distributed shards with privacy-preserving
noise and an aggregation gate. At the scaffold layer this would mean:
per-tenant artifact pools with tenant-tagged provenance, aggregation
across tenants with operator-tunable privacy bounds, an automated eval
gate (held-out task suite per tenant) blocking promotion when patch loss
does not decrease, versioned lineage with RBAC carrying tenant access
policies forward through the promotion, and a rollback envelope on every
promoted artifact. Naming the corner converts a vague gap into a
specific four-property audit criterion; the implementor's safety
requirements are deferred to follow-on practitioner work. It is no
longer ``no system closes both loops'' --- too coarse --- but a specific
four-property combination that the survey can audit any future system
against.

Whether this corner \emph{should} be filled is deployment-dependent. In
safety-critical domains --- medical decision support, legal advice,
financial advice --- reviewer-judgment gates may be a design feature
rather than a defect, and the human latency is the point. The narrower
survey claim is what matters: no public system currently combines
automated local scaffold evolution with governed cross-organisation
promotion, versioned lineage with RBAC, and rollback. Naming the corner
converts a vague gap into a specific four-property audit criterion that
any future system can be measured against. The shortest demonstrable
path appears to build on AlphaEvolve's cascade evaluator + MAP-Elites
archive, extending federation from within-DeepMind multi-target to
cross-customer multi-tenant --- but the architectural surface, not the
timeline, is what this paper claims.

\subsubsection{9.1 Patch-completeness as an operational
hypothesis}\label{patch-completeness-as-an-operational-hypothesis}

The two-tier architecture also raises a measurable question about the
meaning of ``general'' capability in production. We define
\textbf{patch-completeness} as a testable operational condition rather
than a philosophical claim about generality:

\begin{quote}
A model--scaffold composite is \emph{(r, n, d)-patch-complete} for a
deployment patch \(D\) when its residual error rate on a stationary
patch eval suite of size \(n\) stays at or below \(r\) over a deployment
window of \(d\) days.
\end{quote}

This definition makes patch-completeness an empirical property of a
specific deployment and a specific eval, not of a system in general. It
is \emph{not} a claim that any deployed system is patch-complete in the
colloquial sense, nor that a patch-complete system in this sense is
intelligent in any global sense. The values of \(r\), \(n\), and \(d\)
depend on the patch: a coding agent inside one repository, a legal agent
inside one firm's due-diligence workflow, or a support agent on one
product surface each implies its own choice.

We are not aware of any surveyed system that has formally demonstrated
\((r, n, d)\)-patch-completeness for a fixed \(r\), \(n\), \(d\) with
held-out evaluation and a stationary patch over a non-trivial \(d\). The
framework predicts such demonstrations should appear in narrow
production verticals --- a coding agent confined to one repository with
stable conventions, a legal agent confined to a stable due-diligence
workflow --- once eval discipline catches up with deployment scope. The
governance question is then operational: at what residual rate \(r\) and
over what window \(d\) is local competence load-bearing enough to
require governance attention even before global robustness is
established? §10 catalogues the failure surface that the answer would
have to account for.

\subsection{10. A new failure surface}\label{a-new-failure-surface}

Patchwork scaffold maintenance introduces failure modes weight-only
systems do not have; instrumenting the loop to Stage 3 (§2.5) is what
mitigates them at the architectural level, with the specific
countermeasures deferred to follow-on practitioner work. We name the
failure surface here; the full taxonomy (nine categories A--I with
severity ratings and counter-principles P1--P8) is deferred to that
companion.

\textbf{Scaffold bloat} is the dominant near-term risk. IFScale
(Jaroslawicz et al., 2025) reports latency growing roughly 35× for one
reasoning model (o4-mini, 12.4s → 436s) as instruction count scales from
10 to 250, and the best-performing frontier model in that benchmark
(Gemini 2.5 Pro) drops to 68\% accuracy at 500 instructions; weaker
models degrade further (GPT-4o to 15\%). The result is a 2025 snapshot,
and a 2026 follow-up using a later model generation reports the accuracy
degradation softening substantially on the same task family, so the
precise numerical collapse is model-generation-dependent; the underlying
mechanism (lost-in-the-middle, instruction interference; Liu et al.,
2023) persists, and scaffolds that accumulate hundreds of unpruned
instructions remain exposed even when individual frontier-model scaling
improves. \textbf{Poisoned memory} is the dominant security risk:
PoisonedRAG (Zou et al., 2024) reports 90\%+ attack success rate against
standard RAG; Memory Control Flow Attacks (Xu et al., 2026) report
\textgreater90\% vulnerability across major LLM agent frameworks.
\textbf{Prompt injection via tools and MCP} is the dominant cross-vendor
attack surface, with CVE-2025-54136 (``MCPoison'') demonstrating a
public MCP-server compromise pathway. \textbf{Eval fragility} is the
dominant ecosystem risk: the Leaderboard Illusion (Singh et al., 2025)
documents up to 112\% relative performance gains on the Chatbot Arena
distribution under differential training-data access --- a structural
overfitting pathway, not benchmark error; eval suites are themselves
patch-plastic artifacts and must survive Goodhart's Law.
\textbf{Coordination failures} in multi-agent systems are documented in
MAST (Cemri et al., 2025) --- 14 distinct failure modes with
\(\kappa = 0.88\) agreement, dominantly role-confusion / context-loss /
coordination rather than prompt quality. \textbf{Overfitting to patch}
is structural: scaffolds tuned to a deployment over-fit it, and the
framework predicts but does not measure how badly. \textbf{Plasticity
loss at the scaffold level} is the slow-moving risk: context rot,
instruction conflict, memory pollution, and tool-version drift produce a
scaffold-level analogue of weight-level plasticity loss (Dohare et
al.~2024, Lyle et al.~2023). Scaffolds age; a six-month-old
\texttt{CLAUDE.md} is often worse than a clean rebuild, and no surveyed
system treats pruning or expiration as a first-class operation.

Taken together: the patch-plastic surface is real, quantified by
2024--2026 work, and largely unmitigated in production. Follow-on
practitioner work walks the mitigation principles in detail.

\subsection{11. Conclusion}\label{conclusion}

The field has built an external adaptation layer for LLMs, but not the
optimizer that should govern it. If reliability is increasingly repaired
through persistent scaffold artifacts, then production AI needs a theory
and engineering discipline for the discrete, gated, auditable learning
that happens there --- not just for the differentiable learning inside
the model. This paper is one step toward that discipline: a survey of
where production sits on the \emph{patchwork → random-search → descent}
gradient of §2.5, a formalisation of the Stage-3 limit under stated
assumptions, and a named empty corner where the optimizer most plainly
does not yet exist.

The field did not wait for frontier weights to become perfectly
reliable. It built a second learning system around them. That second
system is made of markdown files, skills, memories, tools, orchestration
graphs, eval suites, PRs, version histories, and rollback buttons. It
looks like engineering clutter until viewed as a scaffold parameter
vector under update. Then the clutter resolves into architecture:
production LLMs learn locally, outside the weights --- and most of them
are still learning by patchwork.

The mechanisms were already there. A Claude Skill has been a procedural
weight since Anthropic shipped progressive disclosure. A Cursor Rule has
been a local parameter since \texttt{.cursor/rules/} became a directory.
A PR against \texttt{agents.md} has been a promotion gate since the
Linux Foundation took stewardship. What changed in 2025--2026 is that
the \emph{combination} --- these substrates wrapping a frozen frontier
model and updating from deployment signal --- became dense enough across
vendors that the convergence reads as architecture rather than
coincidence. Once you name it, the empty corners of the design space
become concrete engineering targets. The most consequential one is
enterprise full-auto: Auto-gated Loop 1 plus governed multi-tenant
cross-org Loop 2, with versioned lineage, RBAC, and rollback. In ML
language, a DP-FedAvg-style cross-tenant aggregator with an automated
eval gate at the scaffold layer. Whether it should be filled is
deployment-dependent; that it has not been is the survey's headline
finding.

The five surveyed absences of §6.3 --- failure-triggered memory rare in
production, cross-customer skill promotion absent, memory versioning
unsolved, scaffold-level curriculum unspecified, Loop-2 overfitting
detection missing --- are each a Paper-4 target. None of these is novel
as a complaint; each is novel as a \emph{gap the architecture makes
visible}. If \((r, n, d)\)-patch-complete systems begin to appear in
narrow production verticals --- as §9.1's operational definition makes
possible to test --- then scaffold governance becomes load-bearing
before global robustness is demonstrated.

The deepest connection --- back to Paper 2 --- is structural.
Patch-plasticity is plastic because residual errors are clustered; if
mode discovery were not slow (Postulate 1 in \emph{Architecture of
Errors}), patches would not generalise across deployments and the
architecture would collapse. The survey is consistent with the postulate
at the system level: deployments do generalise enough of the patch to
make scaffold investment pay off. Whether this remains true at agentic,
scientific, and long-horizon scales is the empirical test the next paper
should design.

Practical generality therefore arrives locally first: not measured at
the global model level, but as model--scaffold systems whose competence
is fitted to bounded worlds and is in principle measurable through the
\((r, n, d)\)-patch-completeness condition of §9.1. The frontier model
generalises. The localhost scaffold specialises. Reliability comes from
governing that specialisation. The mechanisms were already there ---
markdown files, skills, memories, tools, orchestration graphs, eval
suites, PRs, version histories, rollback buttons --- but they read as
engineering clutter until viewed as a scaffold parameter vector under
update. Then the clutter resolves into architecture: production LLMs
learn locally, outside the weights, because under deployment-time
constraints that is where patch-specific adaptation can be safely placed
and audited.

\begin{center}\rule{0.5\linewidth}{0.5pt}\end{center}

\subsection{Appendix A. Representative survey
rows}\label{appendix-a.-representative-survey-rows}

The full per-system spreadsheet is \texttt{p3\_master\_scores.csv} in
the repository. The 30 rows below are a representative spot-check
covering all six substrates, both research and production, and all five
clusters of §6.2. \emph{Loop 1} and \emph{Loop 2} columns apply to
systems satisfying the composite-two-loop criterion; \emph{n/a} means
the loop is absent for that system. The \emph{PP} column is the
integrated patch-plasticity score (0--5) used in the survey; the
\emph{Evidence} column abbreviates the strongest available source tier
(T1 peer-reviewed; T2 arXiv preprint; T3 production claim; T4
open-source artifact; T5 industry blog).

\begingroup\small

{\def\LTcaptype{none} % do not increment counter
\begin{longtable}[]{@{}
  >{\raggedright\arraybackslash}p{(\linewidth - 14\tabcolsep) * \real{0.1212}}
  >{\raggedleft\arraybackslash}p{(\linewidth - 14\tabcolsep) * \real{0.0404}}
  >{\raggedright\arraybackslash}p{(\linewidth - 14\tabcolsep) * \real{0.1111}}
  >{\raggedright\arraybackslash}p{(\linewidth - 14\tabcolsep) * \real{0.0707}}
  >{\raggedright\arraybackslash}p{(\linewidth - 14\tabcolsep) * \real{0.0707}}
  >{\raggedleft\arraybackslash}p{(\linewidth - 14\tabcolsep) * \real{0.0404}}
  >{\raggedright\arraybackslash}p{(\linewidth - 14\tabcolsep) * \real{0.1818}}
  >{\raggedright\arraybackslash}p{(\linewidth - 14\tabcolsep) * \real{0.3636}}@{}}
\toprule\noalign{}
\begin{minipage}[b]{\linewidth}\raggedright
System
\end{minipage} & \begin{minipage}[b]{\linewidth}\raggedleft
Year
\end{minipage} & \begin{minipage}[b]{\linewidth}\raggedright
Substrate
\end{minipage} & \begin{minipage}[b]{\linewidth}\raggedright
Loop 1
\end{minipage} & \begin{minipage}[b]{\linewidth}\raggedright
Loop 2
\end{minipage} & \begin{minipage}[b]{\linewidth}\raggedleft
PP
\end{minipage} & \begin{minipage}[b]{\linewidth}\raggedright
Evidence
\end{minipage} & \begin{minipage}[b]{\linewidth}\raggedright
Why included
\end{minipage} \\
\midrule\noalign{}
\endhead
\bottomrule\noalign{}
\endlastfoot
NanoResearch & 2026 & Integrated & Auto & Auto & 5 & T2
(arXiv:2605.10813) & Tri-level Skills/Memory/Policy co-evolution;
cleanest full-auto exemplar \\
AutoAgent & 2026 & Integrated & Auto & Auto & 5 & T2 (arXiv:2603.09716)
& Dual-cycle Execution+Evolution + elastic memory orchestration \\
SkillRL & 2026 & Integrated & Auto & Auto & 5 & T2 (arXiv:2602.08234) &
Recursive skill-augmented RL with \(\text{SR}<0.4\) threshold gate \\
AlphaEvolve & 2025 & Integrated & Auto & Auto & 4 & T2/T3
(arXiv:2506.13131) & Production proof-point: Borg +0.7\%, Gemini kernel
+23\%, FlashAttention +32.5\% \\
DGM (primary) & 2025 & Integrated & Auto & Auto & 4 & T2
(arXiv:2505.22954) & Recursive self-modifying code; SWE-bench
20\%→50\% \\
ADAS & 2024 & Integrated & Auto & Auto & 4 & T2 (arXiv:2408.08435) &
Meta-agent search; Turing-complete agent representation \\
SkillForge & 2026 & Integrated & Auto & Human & 4 & T2
(arXiv:2604.08618) & Production hybrid exemplar; LLM-judge
\textgreater90\% + VFS versioning \\
CASCADE & 2025 & Integrated & Auto & Human & 4.5 & T2 (arXiv:2512.23880)
& Largest published ablation gain in corpus: +57.9 pp on
SciSkillBench \\
EvolveR & 2025 & Integrated & Auto & Auto & 3 & T2 (arXiv:2510.16079) &
Experience-driven lifecycle: episodic memory + skill consolidation +
cross-task transfer; accepted at ICML 2026 \\
Voyager & 2023 & Integrated & Auto & Auto & 4 & T2 (arXiv:2305.16291) &
Foundational skill-library result; 3.3× items, 15.3× tech-tree
milestones \\
AGENTS.md & 2025 & S1 Instructions & n/a & n/a & 2 & T3/T4 (agents.md) &
Cross-vendor standard; 60k+ repos; Lulla 2026 measures −28.6\%
runtime \\
Claude Code CLAUDE.md & 2025 & S1 Instructions & n/a & n/a & 4 & T3/T4
(docs.anthropic.com) & Four scoping levels including org-IT policy +
auto-memory layer; ceiling for S1 \\
Cursor Rules & 2024 & S1 Instructions & n/a & n/a & 3 & T3/T4
(\url{docs.cursor.com/rules}) & \texttt{.cursor/rules/*.mdc} multi-scope
+ MCP attach; git-tracked \\
Anthropic Agent Skills & 2025 & S2 Skills & n/a & Human & 3 & T3
(\url{anthropic.com/engineering}) & Progressive disclosure spec;
LangChain replication 29\%→95\% pass-rate \\
SAGE & 2025 & S2 Skills & Auto & Auto & 5 & T2 (arXiv:2512.17102) &
+8.9\% SGC, 26\% fewer steps, 59\% fewer tokens on AppWorld \\
COSPLAY & 2026 & S2 Skills & Auto & Auto & 5 & T2 (arXiv:2604.20987) &
Co-evolving decision agent + learnable skill bank; skills discovered
from rollouts form a structured skill library for long-horizon tasks \\
Reflexion & 2023 & S3 Memory & Auto & n/a & 4 & T1 (NeurIPS 2023;
arXiv:2303.11366) & Canonical failure-triggered episodic memory; +8\%
HotpotQA \\
Generative Agents & 2023 & S3 Memory & Auto & n/a & 4 & T1 (UIST 2023;
arXiv:2304.03442) & Reflection + importance memory primitive \\
Cuadros et al.~governed collaborative memory & 2026 & S3 Memory & Auto &
Auto & 5 & T2 (arXiv:2605.04264) & Only surveyed memory system with full
provenance + versioning + correction \\
MCP & 2024 & S4 Tools & n/a & n/a & 5 & T3 (modelcontextprotocol.io) &
Cross-vendor standard; thousands of public servers; substantial
enterprise adoption \\
Harvey & 2026 & S4 Tools & n/a & n/a & 5 & T3 (harvey.ai) & 400K
queries/day; 18,000+ workflows; 200+ legal data sources \\
Hippocratic AI & 2026 & S4 Tools & n/a & n/a & 5 & T3
(hippocraticai.com) & \$3.5B valuation; 30\% readmission reduction;
360\% care-capacity boost \\
Multi-Agent Evolve (MAE) & 2025 & S5 Orchestration & Auto & n/a & 4 & T2
(arXiv:2510.23595) & RL co-evolution of Proposer/Solver/Judge
population; one loop closed, no Loop-2 cross-context promotion \\
Cemri et al.~(AG2) & 2025 & S5 Orchestration & n/a & n/a & 3 & T2
(arXiv:2503.13657) & Specialisation +4.5 pp on GSM-Plus (\(p = 0.03\));
MAST 14 failure modes \\
Braintrust & 2024 & S6 Governance & n/a & Auto & 5 & T3 (braintrust.dev)
& Prompts-as-code: immutable commits, eval-gated PR merge, sub-5 min
rollback \\
Vellum & 2024 & S6 Governance & n/a & Auto & 5 & T3 (vellum.ai) & Same
governance pattern; production deployment with eval thresholds \\
LangSmith Hub & 2024 & S6 Governance & n/a & Auto & 5 & T3
(\url{smith.langchain.com}) & Prompt repository with environment
promotion and eval blocking \\
Self-Refine* & 2023 & (contrast) & n/a & n/a & 1 & T2 (arXiv:2303.17651)
& Fails the patch-plastic discriminator: in-context iteration only \\
Mem0* & 2024 & (contrast) & n/a & n/a & 2 & T3/T4 (mem0.ai) & Fails Loop
2: per-user memory with no cross-context promotion path \\
Tran \& Kiela* & 2026 & (contrast) & n/a & n/a & 1 & T2
(arXiv:2604.02460) & Single-agent baselines match or beat multi-agent
systems under matched thinking-token budgets --- orthogonal finding
constraining S5 claims \\
\end{longtable}
}

\endgroup

\emph{Asterisks mark contrast systems included for definitional
clarity.}

\begin{center}\rule{0.5\linewidth}{0.5pt}\end{center}

\subsection{Appendix B. Data appendix}\label{appendix-b.-data-appendix}

\textbf{Corpus size.} The supplementary file
\texttt{p3\_master\_scores.csv} contains approximately 130 unique LLM
systems scored across the six substrates. Because a system can be scored
on multiple substrates, the row count in the CSV exceeds the
unique-system count; the system count cited throughout §§5--8 refers to
unique systems.

\textbf{Core / contrast split.} Approximately 90 unique systems satisfy
the patch-plastic discriminator (score ≥ 3 on at least one substrate)
and form the \textbf{core corpus}. The remaining ≈ 38 systems form the
\textbf{contrast corpus} --- commonly described as agentic or
self-improving in public discourse but failing at least one
discriminator. Headline statistics quote the core corpus; the contrast
set is preserved with asterisks for definitional clarity. Appendix D
enumerates the contrast set.

\textbf{Evidence-tier distribution.} Of the core corpus, approximately
one quarter carry T1 (peer-reviewed) evidence, roughly half carry T2
(arXiv preprint) evidence, and the remainder carry T3 (production claim
with metrics) or T4 (open-source artifact) evidence as the primary tier.
T5 (industry-blog/interview) is admissible only when corroborated by a
higher tier elsewhere. The per-row evidence flag in Appendix A and in
the CSV makes the propagation auditable.

\textbf{Sensitivity analysis (composite two-loop count).} The strict
count of thirteen research systems satisfying the composite-two-loop
criterion changes under two natural relaxations of the gate definition:
under (a) --- counting any of versioning, lineage, or rollback as Loop-2
governance --- the count rises to approximately seventeen by
re-admitting MAE, MetaGen, MetaReflection, and CLIN; under (b) ---
admitting session signal that updates optimizer state without persisting
an artifact distinct from optimiser state --- the count rises to
approximately fifteen by re-admitting DSPy and TextGrad. The qualitative
findings (Auto×Auto well-populated, {[}Human × Human{]} and {[}Human ×
Auto{]} empty in the surveyed corpus, governance-first cluster in
{[}None × Auto{]}, enterprise full-auto cell empty) survive both
relaxations.

\textbf{Scoring discipline.} Scoring is single-rater. The published 0--5
rubric in §4 is the canonical scoring criterion; half-scores were used
sparingly when one criterion was clearly met and another partially.
Ambiguous cases were resolved conservatively (downward). The CSV exposes
every per-system score so reviewers can audit ambiguous classifications
independently against the per-substrate evidence harvests
(\texttt{p3\_harvest\_s1\_instructions.md} through
\texttt{p3\_harvest\_s6\_governance.md},
\texttt{p3\_harvest\_integrated.md}).

\textbf{Survey window.} January 2023 -- May 2026. Preprint-only systems
carry T2 evidence and are flagged in Appendix A. A fraction of T2 work
is recent enough that ablations may yet be revised; cluster-level claims
are robust to dropping any single T2 system.

\textbf{Supplementary files (shipped in the arXiv source bundle under
\texttt{supplementary/}).}

{\def\LTcaptype{none} % do not increment counter
\begin{longtable}[]{@{}
  >{\raggedright\arraybackslash}p{(\linewidth - 2\tabcolsep) * \real{0.5000}}
  >{\raggedright\arraybackslash}p{(\linewidth - 2\tabcolsep) * \real{0.5000}}@{}}
\toprule\noalign{}
\begin{minipage}[b]{\linewidth}\raggedright
File
\end{minipage} & \begin{minipage}[b]{\linewidth}\raggedright
Contents
\end{minipage} \\
\midrule\noalign{}
\endhead
\bottomrule\noalign{}
\endlastfoot
\texttt{p3\_master\_scores.csv} & Per-system, per-substrate 0--5 scores
with evidence citations \\
\texttt{p3\_harvest\_s1\_instructions.md} & Evidence harvest for S1
(instructions) \\
\texttt{p3\_harvest\_s2\_skills.md} & Evidence harvest for S2
(skills) \\
\texttt{p3\_harvest\_s3\_memory.md} & Evidence harvest for S3
(memory) \\
\texttt{p3\_harvest\_s4\_tools.md} & Evidence harvest for S4 (tools) \\
\texttt{p3\_harvest\_s5\_orchestration.md} & Evidence harvest for S5
(orchestration) \\
\texttt{p3\_harvest\_s6\_governance.md} & Evidence harvest for S6
(governance) \\
\texttt{p3\_harvest\_integrated.md} & Cross-substrate integrated
audit \\
\texttt{p3\_harvest\_composite\_two\_loop.md} & Strict
composite-two-loop audit, 22 candidates \\
\texttt{p3\_harvest\_failure\_modes.md} & Evidence harvest for §10
failure surface \\
\texttt{p3\_prior\_art\_check.md} & Prior-art reconciliation against
adjacent literatures \\
\end{longtable}
}

\begin{center}\rule{0.5\linewidth}{0.5pt}\end{center}

\subsection{Appendix C. Gate type → estimator-family
mapping}\label{appendix-c.-gate-type-estimator-family-mapping}

The §3.2 formalism collapses production gates into a single per-edit
gate \(G_t\) that estimates the directional loss change
\(\widehat{\Delta}_{z_t} L_D\). In practice, surveyed gates instantiate
that abstraction in several distinct ways, and some are better read as
\emph{proposal-pruning} steps that precede the per-edit gate rather than
as direct \(\widehat{\Delta}\) estimators. The table below maps each
surveyed gate family to its role in the formalism.

{\def\LTcaptype{none} % do not increment counter
\begin{longtable}[]{@{}
  >{\raggedright\arraybackslash}p{(\linewidth - 6\tabcolsep) * \real{0.2500}}
  >{\raggedright\arraybackslash}p{(\linewidth - 6\tabcolsep) * \real{0.2500}}
  >{\raggedright\arraybackslash}p{(\linewidth - 6\tabcolsep) * \real{0.2500}}
  >{\raggedright\arraybackslash}p{(\linewidth - 6\tabcolsep) * \real{0.2500}}@{}}
\toprule\noalign{}
\begin{minipage}[b]{\linewidth}\raggedright
Gate family
\end{minipage} & \begin{minipage}[b]{\linewidth}\raggedright
What it estimates
\end{minipage} & \begin{minipage}[b]{\linewidth}\raggedright
Slot in §3.2
\end{minipage} & \begin{minipage}[b]{\linewidth}\raggedright
Representative systems
\end{minipage} \\
\midrule\noalign{}
\endhead
\bottomrule\noalign{}
\endlastfoot
\textbf{Eval-cascade} & \(\widehat{\Delta} L_D\) on a held-out task
suite; cascaded through cheap-to-expensive evaluators & Per-edit gate
\(G_t\) & AlphaEvolve, Braintrust, Vellum, LangSmith Hub \\
\textbf{RL-threshold} & Reward improvement vs.~a deployment threshold
(e.g., task-category accuracy \textless{} 0.4 triggers replacement) &
Per-edit gate \(G_t\) & SkillRL, SAGE, EvolveR \\
\textbf{LLM-judge with calibration} & Pairwise- or rubric-scored
quality, calibrated against human agreement (LLM-judge ≥ 90\% agreement,
etc.) & Per-edit gate \(G_t\) & SkillForge (Loop 1), CASCADE
(consolidation gate) \\
\textbf{Surrogate verifier} & Symbolic, learned, or test-case verifier
producing a pass/fail or graded score & Per-edit gate \(G_t\) & DGM
(test-pass on sandboxed code), Voyager (in-game verifier) \\
\textbf{Reviewer-judgment} & Human credit assignment via PR review or
expert sign-off; no algorithmic threshold & Per-edit gate \(G_t\)
(human-instantiated) & SkillForge (Loop 2), Sierra Agent OS 2.0,
Anthropic skills repo \\
\textbf{Semantic-merge} & Whether a candidate is semantically duplicated
by, or compatible with, existing scaffold artifacts &
\emph{Proposal-pruning} preceding \(G_t\) & NanoResearch (orchestrator
semantic-merge step) \\
\textbf{MAP-Elites / archive selection} & Whether a candidate populates
a previously unfilled cell in a behavior-diversity archive &
\emph{Proposal-pruning} preceding \(G_t\) & AlphaEvolve (archive), ADAS
(meta-agent search) \\
\textbf{Multi-sample agreement} & Whether \(k\) independent proposals or
evaluations converge on the same direction; false-accept rate decays as
\(p^k\) rather than variance shrinking as \(1/k\) & \emph{Aggregation
filter} (qualitatively distinct from mini-batch averaging) &
Under-explored in the surveyed corpus \\
\end{longtable}
}

The eval-cascade, RL-threshold, LLM-judge, surrogate-verifier, and
reviewer-judgment families directly populate the per-edit gate \(G_t\)
of §3.2 and inherit its descent guarantees (under Assumptions A1--A5).
The semantic-merge and archive-selection families are not estimators of
\(\widehat{\Delta} L_D\) --- they constrain the proposal distribution
\(Q_t\) rather than judging individual candidate quality, and the
descent argument therefore applies to the composition (proposal-pruning
∘ per-edit gate) rather than to either step alone. Multi-sample
agreement filters are not surveyed in production form; they are listed
for completeness because they are the correct ML analogue for ``wait for
\(k\) pieces of evidence before accepting an edit,'' which is \emph{not}
the same operation as mini-batch averaging and should not be conflated
with it in scaffold-descent analyses.

\begin{center}\rule{0.5\linewidth}{0.5pt}\end{center}

\subsection{Appendix D. Contrast set}\label{appendix-d.-contrast-set}

The contrast corpus of approximately 38 systems is included to make the
patch-plastic discriminator audit-able. Each system here is commonly
described in public discourse as agentic, self-improving, or scaffolded,
but fails at least one discriminator. Representative entries:

{\def\LTcaptype{none} % do not increment counter
\begin{longtable}[]{@{}
  >{\raggedright\arraybackslash}p{(\linewidth - 6\tabcolsep) * \real{0.2500}}
  >{\raggedright\arraybackslash}p{(\linewidth - 6\tabcolsep) * \real{0.2500}}
  >{\raggedright\arraybackslash}p{(\linewidth - 6\tabcolsep) * \real{0.2500}}
  >{\raggedright\arraybackslash}p{(\linewidth - 6\tabcolsep) * \real{0.2500}}@{}}
\toprule\noalign{}
\begin{minipage}[b]{\linewidth}\raggedright
System
\end{minipage} & \begin{minipage}[b]{\linewidth}\raggedright
Year
\end{minipage} & \begin{minipage}[b]{\linewidth}\raggedright
Primary discriminator failure
\end{minipage} & \begin{minipage}[b]{\linewidth}\raggedright
Brief reason
\end{minipage} \\
\midrule\noalign{}
\endhead
\bottomrule\noalign{}
\endlastfoot
Self-Refine (Madaan et al.) & 2023 & No persistent artifact & In-context
iteration only; no scaffold delta survives the session \\
ReAct (Yao et al.) & 2023 & No persistent artifact & Thought-act-observe
loop is intra-session; no Loop 1 persistence \\
Reflexion (Shinn et al.) & 2023 & No Loop 2 (Loop 1 only) &
Failure-triggered episodic memory persists per-agent but no
cross-context promotion mechanism \\
Generative Agents (Park et al.) & 2023 & No Loop 2 & Reflection +
importance memory primitive; per-agent scope \\
AutoManual (Chen et al.) & 2024 & Passes (in core corpus) & Listed here
only to flag boundary case; see Appendix A \\
ExpeL & 2023 & Loop boundary blurred & Experience bank built from
training tasks, not in-deployment failures \\
Mem0 (vanilla deployment) & 2024 & No Loop 2 & Per-user memory store
with no cross-context promotion path \\
Zep & 2024 & No Loop 2 & Temporal knowledge graph per agent/session, no
promotion gate \\
ChatGPT Memory & 2024 & No Loop 2 & Per-user persistence; no
inter-tenant promotion \\
Claude Memory & 2025 & No Loop 2 & Per-conversation-thread persistence;
no inter-tenant promotion \\
MemGPT & 2023 & No Loop 2 & Virtual-context management per agent; no
promotion mechanism \\
DSPy & 2023 & No artifact distinct from optimiser state &
Prompt-template optimization but optimised templates are optimiser
state, not persisted scaffold \\
TextGrad & 2024 & No artifact distinct from optimiser state & Same
boundary case as DSPy \\
MetaGen & 2026 & No Loop 2; no cross-session persistence & Role pool
dynamic within single inference call \\
MAE (Multi-Agent Evolve) & 2025 & No Loop 2 & Co-evolution within
Proposer/Solver/Judge triplet; no external promotion \\
MetaReflection & 2024 & No Loop 2 & Per-agent semantic memory; no
cross-agent promotion \\
CLIN & 2023 & No cross-instance promotion & Causal abstractions
per-agent-per-environment \\
LangGraph & 2024 & Framework only & Orchestration framework; not a
deployed patch-plastic system \\
AutoGen & 2023 & Framework only & Multi-agent framework; not a deployed
patch-plastic system \\
CrewAI & 2024 & Framework only & Multi-agent framework; not a deployed
patch-plastic system \\
Tran \& Kiela & 2026 & Single-agent baseline & Listed because the result
constrains S5 maturity claims (see §7) \\
W\&B Weave & 2024 & No Loop 1 & Observability/eval-as-code;
engineer-authored artifacts \\
Agenta & 2024 & No Loop 1 & Prompt-as-code platform; engineer-authored
artifacts \\
LangFuse & 2024 & No Loop 1 & Observability/tracing; engineer-authored
artifacts \\
Mobile-Agent-E & 2025 & Loop 2 absent & Self-evolving mobile assistant;
no cross-tenant promotion \\
Cradle & 2024 & Loop 2 absent & Foundation agent for general computer
control; per-session only \\
OS-Copilot & 2024 & Loop 2 absent & Computer-use generalist;
per-deployment \\
Agent S, Agent S2 & 2024--2025 & Loop 2 absent & Computer-use agentic
framework; no cross-context promotion \\
ACE (Zhang et al.) & 2025 & No Loop 2 in surveyed form & Single-context
evolving playbook; per-deployment \\
Claude Code memory (vanilla CLAUDE.md) & 2025 & No Loop 2
(within-project only) & Lives in {[}None × None{]} without the cross-org
governance layer \\
Cursor Rules (vanilla) & 2024 & No Loop 2 & Per-project artifacts; no
cross-org promotion mechanism \\
GitHub Copilot custom instructions & 2024 & No Loop 1 &
Engineer-authored at repo base branch \\
Windsurf rules & 2024 & No Loop 1 & Engineer-authored, no
failure-triggered update \\
Continue.dev rules & 2024 & No Loop 1 & Same as above \\
Zed AI rules & 2024 & No Loop 1 & Same as above \\
Replit \texttt{replit.md} & 2025 & Loop 2 absent & Self-writing scaffold
but per-workspace only \\
Aider conventions & 2024 & No Loop 1 & Engineer-authored convention
file \\
Lovable knowledge & 2024 & No Loop 1 & Engineer-authored knowledge
field \\
\end{longtable}
}

The full per-system exclusion log lives in
\texttt{supplementary/p3\_master\_scores.csv} (with one row per
substrate scoring) and in the per-substrate evidence harvests. The list
above is representative, not exhaustive; some contrast entries are
excluded from the table for length and appear only in the CSV.

\begin{center}\rule{0.5\linewidth}{0.5pt}\end{center}

\subsection{References}\label{references}

AGENTS.md. (2025). AGENTS.md: A cross-vendor open standard for agent
instructions. Linux Foundation. \url{https://agents.md}

Agashe, S., Han, J., Gan, S., Yang, J., Li, A., \& Wang, X. E. (2024).
Agent S: An open agentic framework that uses computers like a human.
\emph{arXiv preprint arXiv:2410.08164}.

Agashe, S., et al.~(2025). Agent S2: A compositional
generalist-specialist framework for computer use agents. \emph{arXiv
preprint arXiv:2504.00906}.

Anthropic. (2024). Model Context Protocol specification.
\url{https://modelcontextprotocol.io}

Arbuzov, M. L., Bei, S., Dong, Z., Kalaev, D., \& Shvets, A. A. (2025).
Beyond exponential decay: Rethinking error accumulation in large
language models. \emph{arXiv preprint arXiv:2505.24187}.

Arbuzov, M. L., Shvets, A. A., \& Bei, S. (2026). The architecture of
errors: Logarithmic mode discovery and polylogarithmic intervention
budgets for long-context LLM reliability. \emph{arXiv preprint}
(forthcoming).

Cemri, M., et al.~(2025). Why do multi-agent LLM systems fail?
\emph{arXiv preprint arXiv:2503.13657}. Introduces the MAST taxonomy.

Chen, M., Li, Y., Yang, Y., Yu, S., Lin, B., \& He, X. (2024).
AutoManual: Constructing instruction manuals by LLM agents via
interactive environmental learning. \emph{NeurIPS 2024}.
arXiv:2405.16247.

Chen, Y., et al.~(2025). Multi-Agent Evolve: LLM self-improve through
co-evolution. \emph{arXiv preprint arXiv:2510.23595}.

Cuadros, D. F., Maiga, A.-A., Meskhidze, H., \& Curtis-Trudel, A.
(2026). Governed collaborative memory as artificial selection in
LLM-based multi-agent systems. \emph{arXiv preprint arXiv:2605.04264}.

Dohare, S., et al.~(2024). Loss of plasticity in deep continual
learning. \emph{Nature}, 632(8026), 768--774.

Du, P. (2026). Memory for autonomous LLM agents: Mechanisms, evaluation,
and emerging frontiers. \emph{arXiv preprint arXiv:2603.07670}.

Fang, J., Peng, Y., Zhang, X., et al.~(2025). A comprehensive survey of
self-evolving AI agents: A new paradigm bridging foundation models and
lifelong agentic systems. \emph{arXiv preprint arXiv:2508.07407}.

Fawzi, A., Balog, M., Huang, A., et al.~(2022). Discovering faster
matrix multiplication algorithms with reinforcement learning.
\emph{Nature}, 610(7930), 47--53. (DeepMind AlphaTensor.)

Gao, H.-a., Geng, J., Hua, W., et al.~(2026). A survey of self-evolving
agents: What, when, how, and where to evolve on the path to artificial
super intelligence. \emph{Transactions on Machine Learning Research}.
arXiv:2507.21046.

Harada, K., et al.~(2025). Curse of instructions: Large language models
cannot follow multiple instructions at once. \emph{ICLR 2025}.
\url{https://openreview.net/forum?id=R6q67CDBCH}

Harvey. (2026). Harvey product overview. \url{https://www.harvey.ai}

Hippocratic AI. (2026). Polaris clinical outcome evidence.
\url{https://www.hippocraticai.com}

Hu, S., Lu, C., \& Clune, J. (2024). Automated design of agentic systems
(ADAS). \emph{arXiv preprint arXiv:2408.08435}.

Huang, X., et al.~(2025). CASCADE: Cumulative agentic skill creation
through autonomous development and evolution. \emph{arXiv preprint
arXiv:2512.23880}.

Jaroslawicz, D., et al.~(2025). How many instructions can LLMs follow at
once? \emph{arXiv preprint arXiv:2507.11538}. Introduces the IFScale
benchmark.

Jiang, P., Lin, J., Shi, Z., et al.~(2025). Adaptation of agentic AI: A
survey of post-training, memory, and skills. \emph{arXiv preprint
arXiv:2512.16301}.

Kirk, R., Mediratta, I., Nalmpantis, C., et al.~(2024). Understanding
the effects of RLHF on LLM generalisation and diversity. \emph{ICLR
2024}. arXiv:2310.06452.

Liu, N. F., et al.~(2023). Lost in the middle: How language models use
long contexts. \emph{TACL}. arXiv:2307.03172.

Liu, X., et al.~(2026). SkillForge: Forging domain-specific,
self-evolving agent skills in cloud technical support. \emph{arXiv
preprint arXiv:2604.08618}. Accepted at ACM SIGIR 2026 Industry Track.

Lulla, J. L., et al.~(2026). On the impact of AGENTS.md files on the
efficiency of AI coding agents. \emph{arXiv preprint arXiv:2601.20404}.

Lyle, C., Zheng, Z., Nikishin, E., Avila Pires, B., Pascanu, R., \&
Dabney, W. (2023). Understanding plasticity in neural networks.
\emph{ICML 2023} (oral). arXiv:2303.01486.

Madaan, A., et al.~(2023). Self-Refine: Iterative refinement with
self-feedback. \emph{NeurIPS 2023}. arXiv:2303.17651.

Newell, A. (1990). \emph{Unified theories of cognition}. Harvard
University Press.

Novikov, A., et al.~(2025). AlphaEvolve: A coding agent for scientific
and algorithmic discovery. \emph{arXiv preprint arXiv:2506.13131}.
(DeepMind.)

Padmakumar, V., \& He, H. (2023). Does writing with language models
reduce content diversity? \emph{arXiv preprint arXiv:2309.05196}.

Park, J. S., et al.~(2023). Generative agents: Interactive simulacra of
human behavior. \emph{UIST 2023}. arXiv:2304.03442.

Press, O., et al.~(2022). Measuring and narrowing the compositionality
gap in language models. \emph{arXiv preprint arXiv:2210.03350}.

Shinn, N., et al.~(2023). Reflexion: Language agents with verbal
reinforcement learning. \emph{NeurIPS 2023}. arXiv:2303.11366.

Sierra. (2024). Sierra Agent OS. \url{https://sierra.ai}

Sierra. (2025). Agent OS 2.0: From answers to memory and action.
\url{https://sierra.ai/blog/agent-os-2-0}

Singh, S., et al.~(2025). The leaderboard illusion. \emph{arXiv preprint
arXiv:2504.20879}.

Sumers, T. R., Yao, S., Narasimhan, K., \& Griffiths, T. L. (2023).
Cognitive architectures for language agents (CoALA). \emph{Transactions
on Machine Learning Research}. arXiv:2309.02427.

Tan, W., et al.~(2024). Cradle: Empowering foundation agents towards
general computer control. \emph{arXiv preprint arXiv:2403.03186}.

Tran, D., \& Kiela, D. (2026). Single-agent LLMs outperform multi-agent
systems on multi-hop reasoning under equal thinking token budgets.
\emph{arXiv preprint arXiv:2604.02460}.

Wang, G., et al.~(2023). Voyager: An open-ended embodied agent with
large language models. \emph{arXiv preprint arXiv:2305.16291}.

Wang, J., et al.~(2025). Reinforcement learning for self-improving agent
with skill library. \emph{arXiv preprint arXiv:2512.17102}. Method named
SAGE (Skill Augmented GRPO for self-Evolution).

Wang, X., et al.~(2026). AutoAgent: Evolving cognition and elastic
memory orchestration for adaptive agents. \emph{arXiv preprint
arXiv:2603.09716}.

Wang, Z., et al.~(2025). Mobile-Agent-E: Self-evolving mobile assistant
for complex tasks. \emph{arXiv preprint arXiv:2501.11733}.

Wu, R., et al.~(2025). EvolveR: Self-evolving LLM agents through an
experience-driven lifecycle. \emph{arXiv preprint arXiv:2510.16079}.
Accepted at ICML 2026.

Wu, X., et al.~(2026). Co-evolving LLM decision and skill bank agents
for long-horizon tasks. \emph{arXiv preprint arXiv:2604.20987}.
Framework named COSPLAY.

Wu, Z., et al.~(2024). OS-Copilot: Towards generalist computer agents
with self-improvement. \emph{arXiv preprint arXiv:2402.07456}.

Xia, P., et al.~(2026). SkillRL: Evolving agents via recursive
skill-augmented reinforcement learning. \emph{arXiv preprint
arXiv:2602.08234}.

Xu, Jinhang, et al.~(2026). NanoResearch: Co-evolving skills, memory,
and policy for personalized research automation. \emph{arXiv preprint
arXiv:2605.10813}.

Xu, Zhenlin, et al.~(2026). From storage to steering: Memory control
flow attacks on LLM agents. \emph{arXiv preprint arXiv:2603.15125}.
Introduces the MCFA attack class.

Zhang, H., et al.~(2026). CoEvoSkills: Self-evolving agent skills via
co-evolutionary verification. \emph{arXiv preprint arXiv:2604.01687}.

Zhang, J., et al.~(2025). Darwin Gödel machine: Open-ended evolution of
self-improving agents. \emph{arXiv preprint arXiv:2505.22954}. (Body
text refers to this as DGM.)

Zhang, Q., Hu, C., Upasani, S., et al.~(2025). Agentic Context
Engineering: Evolving contexts for self-improving language models.
\emph{arXiv preprint arXiv:2510.04618}.

Zhou, Y., Shu, W., Su, Y., et al.~(2026). A comprehensive survey on
agent skills: Taxonomy, techniques, and applications. \emph{arXiv
preprint arXiv:2605.07358}.

Zou, W., et al.~(2024). PoisonedRAG: Knowledge corruption attacks to
retrieval-augmented generation of large language models. \emph{arXiv
preprint arXiv:2402.07867}. Accepted at USENIX Security 2025.

\end{document}
